% !TeX program = pdflatex
% !BIB program = biber
% ==============================================================================
% Modelo de Trabalho de Conclusão de Curso (TCC) do tipo TECNOLÓGICO
% Ciência da Computação e Engenharia de Computação | DCE/Ceunes/Ufes
%
% Normas: ABNT NBR 14724:2024 (trabalhos acadêmicos), NBR 10520:2023
% (citações), NBR 6023 (referências), NBR 6024:2012 (numeração progressiva),
% NBR 6027:2012 (sumário) e NBR 6028:2021 (resumo).
%
% Adaptado do modelo do DI/CT/Ufes (Prof. Vítor E. Silva Souza), baseado no
% abnTeX2. Distribuído sob a LaTeX Project Public License 1.3c.
%
% Comece pelo arquivo tecnologico/dados.tex. O texto de cada capítulo traz um
% exemplo fictício (um sistema web para agendamento de monitorias)
% e caixas de orientação que explicam o que escrever em cada parte.
% ==============================================================================

\documentclass[
  12pt,           % tamanho da fonte (NBR 14724 recomenda 12)
  a4paper,        % papel A4
  oneside,        % entrega em PDF: cada seção primária começa em nova página
  openany,
  chapter=TITLE,  % títulos das seções primárias em caixa alta
  english,        % idioma adicional (abstract)
  brazil,         % idioma principal (deve ser o último)
]{abntex2}

% ------------------------------------------------------------------------------
% Chaves de controle
% ------------------------------------------------------------------------------
% Caixas de orientação: troque para \mostrarinstrucoesfalse quando não
% precisar mais delas (obrigatório na versão entregue à banca).
\newif\ifmostrarinstrucoes
\mostrarinstrucoestrue

% Notas de revisão \notaaluno{...} e \notaorientador{...}.
\newif\ifmostrarnotas
\mostrarnotastrue

% ------------------------------------------------------------------------------
% Preâmbulo
% ------------------------------------------------------------------------------
% ==============================================================================
% config/pacotes.tex
% Pacotes usados pelo modelo. Em geral, não é preciso alterar este arquivo.
% Se precisar de um pacote novo, acrescente-o ao final, na área indicada.
% ==============================================================================

% ------------------------------------------------------------------------------
% Fontes
% ------------------------------------------------------------------------------
% O modelo compila com pdfLaTeX (padrão), XeLaTeX ou LuaLaTeX. O pacote iftex
% identifica o compilador e carrega apenas o que cada um precisa.
\usepackage{iftex}
\ifPDFTeX
  \usepackage[T1]{fontenc}
  \usepackage{lmodern}   % Latin Modern Mono para códigos (texto e títulos usam Libertine)
\fi

% Linux Libertine (texto) e Linux Biolinum (sem serifa). A opção mono=false
% mantém a fonte monoespaçada padrão, mais legível em listagens de código.
\usepackage[mono=false]{libertine}

% Matemática com o mesmo desenho da Libertine. O libertinust1math já define os
% símbolos do amssymb, por isso este último não deve ser carregado.
\usepackage{amsmath}
\usepackage{libertinust1math}

% Vírgula decimal em modo matemático sem espaço extra: $3,14$.
\usepackage{icomma}

% Ajustes finos de justificação e hifenização.
\usepackage{microtype}

% Recuo também no primeiro parágrafo de cada seção.
\usepackage{indentfirst}

% ------------------------------------------------------------------------------
% Cores, figuras e gráficos
% ------------------------------------------------------------------------------
\usepackage[dvipsnames,table]{xcolor}
\usepackage{graphicx}
\usepackage{subcaption}
\usepackage{pgfplots}
\pgfplotsset{compat=1.18}
\usetikzlibrary{positioning,arrows.meta,shapes.multipart,shapes.geometric,fit,calc,backgrounds}

% ------------------------------------------------------------------------------
% Tabelas
% ------------------------------------------------------------------------------
\usepackage{array}
\usepackage{booktabs}
\usepackage{multirow}
\usepackage{tabularx}
\usepackage{longtable}

% ------------------------------------------------------------------------------
% Códigos e algoritmos
% ------------------------------------------------------------------------------
\usepackage{listings}
\usepackage{algpseudocode}

% ------------------------------------------------------------------------------
% Caixas de orientação e exemplos (podem ser desligadas no arquivo principal)
% ------------------------------------------------------------------------------
\usepackage{comment}
\usepackage[most]{tcolorbox}

% ------------------------------------------------------------------------------
% Citações e referências (ABNT NBR 10520:2023 e NBR 6023)
% ------------------------------------------------------------------------------
% Requer biblatex-abnt 4.0 ou mais recente (TeX Live 2024 em diante) e o Biber
% como processador de bibliografia.
\usepackage[autostyle]{csquotes}
\usepackage[style=abnt, backend=biber]{biblatex}

% ------------------------------------------------------------------------------
% Referências cruzadas e siglas
% ------------------------------------------------------------------------------
\usepackage[nameinlink,noabbrev]{cleveref}
\usepackage{acro}

% ------------------------------------------------------------------------------
% Notas de revisão entre aluno e orientador
% ------------------------------------------------------------------------------
\ifmostrarnotas
  \usepackage[colorinlistoftodos,textsize=footnotesize]{todonotes}
\else
  \usepackage[disable]{todonotes}
\fi

\usepackage{xspace}

% ------------------------------------------------------------------------------
% Pacotes adicionais do aluno
% ------------------------------------------------------------------------------
% \usepackage{...}

% ==============================================================================
% config/estilo.tex
% Aparência do documento: títulos, legendas, ilustrações, códigos, referências
% cruzadas, siglas e caixas de orientação. Em geral, não é preciso alterar.
% ==============================================================================

% ------------------------------------------------------------------------------
% Termo usado para as divisões principais do texto
% ------------------------------------------------------------------------------
% A NBR 14724:2024 (item 4.2.2) determina que o trabalho seja organizado em
% seções, e não em capítulos. O padrão abaixo segue as regras de TCC do curso.
% Para adotar a terminologia da norma, troque por {Seção} e {Seções} e revise
% os artigos no texto ("o Capítulo 2" passa a ser "a Seção 2").
\newcommand{\termocap}{Capítulo}
\newcommand{\termocaps}{Capítulos}

% ------------------------------------------------------------------------------
% Parágrafos e títulos
% ------------------------------------------------------------------------------
\setlength{\parindent}{1.25cm}

% Destaque gráfico progressivo das seções (NBR 6024). O caixa alta das seções
% primárias vem da opção chapter=TITLE da classe.
\renewcommand{\ABNTEXchapterfont}{\normalfont\bfseries}
\renewcommand{\ABNTEXchapterfontsize}{\large}
\renewcommand{\ABNTEXsectionfont}{\normalfont\bfseries}
\renewcommand{\ABNTEXsectionfontsize}{\normalsize}
\renewcommand{\ABNTEXsubsectionfont}{\normalfont\bfseries\itshape}
\renewcommand{\ABNTEXsubsectionfontsize}{\normalsize}
\renewcommand{\ABNTEXsubsubsectionfont}{\normalfont\itshape}
\renewcommand{\ABNTEXsubsubsectionfontsize}{\normalsize}

% ------------------------------------------------------------------------------
% Legendas e fontes das ilustrações e tabelas
% ------------------------------------------------------------------------------
% A NBR 14724 pede fonte menor e uniforme para legendas, fontes, notas e
% citações longas. O abnTeX2 usa \footnotesize para \fonte e citações longas.
\captionsetup{font=footnotesize}
\captionsetup[sub]{font=footnotesize}

% Texto padrão para material elaborado pelo próprio autor (NBR 14724:2024).
\newcommand{\fonteautor}{\fonte{Elaborado pelo próprio autor.}}

% Posição padrão dos objetos flutuantes.
\setfloatlocations{figure}{htbp}
\setfloatlocations{table}{htbp}

% ------------------------------------------------------------------------------
% Novos tipos de ilustração: Quadro, Código e Algoritmo
% ------------------------------------------------------------------------------
% Quadro: informação textual organizada em linhas e colunas, com bordas
% fechadas. Tabela: dados numéricos, sem bordas laterais (normas do IBGE).
\newcommand{\quadroname}{Quadro}
\newcommand{\listofquadrosname}{Lista de quadros}
\newfloat{quadro}{loq}{\quadroname}
\newlistof{listofquadros}{loq}{\listofquadrosname}
\newlistentry{quadro}{loq}{0}
\setfloatlocations{quadro}{htbp}
\renewcommand{\cftquadroname}{\quadroname\space}
\renewcommand*{\cftquadroaftersnum}{\hfill\textendash\hfill}

\newcommand{\codigoname}{Código}
\newcommand{\listofcodigosname}{Lista de códigos}
\newfloat{codigo}{lcod}{\codigoname}
\newlistof{listofcodigos}{lcod}{\listofcodigosname}
\newlistentry{codigo}{lcod}{0}
\setfloatlocations{codigo}{htbp}
\renewcommand{\cftcodigoname}{\codigoname\space}
\renewcommand*{\cftcodigoaftersnum}{\hfill\textendash\hfill}

\newcommand{\algoritmoname}{Algoritmo}
\newcommand{\listofalgoritmosname}{Lista de algoritmos}
\newfloat{algoritmo}{lalg}{\algoritmoname}
\newlistof{listofalgoritmos}{lalg}{\listofalgoritmosname}
\newlistentry{algoritmo}{lalg}{0}
\setfloatlocations{algoritmo}{htbp}
\renewcommand{\cftalgoritmoname}{\algoritmoname\space}
\renewcommand*{\cftalgoritmoaftersnum}{\hfill\textendash\hfill}

% Largura reservada para o rótulo ("Quadro 10", "Código 10" etc.) nas listas.
\setlength{\cftquadronumwidth}{2.5em}
\setlength{\cftcodigonumwidth}{2.5em}
\setlength{\cftalgoritmonumwidth}{2.5em}

% ------------------------------------------------------------------------------
% Códigos (pacote listings)
% ------------------------------------------------------------------------------
\definecolor{codigofundo}{gray}{0.96}
\definecolor{codigocomentario}{gray}{0.40}
\definecolor{codigonumero}{gray}{0.55}

\lstset{
  basicstyle=\ttfamily\footnotesize\linespread{1}\selectfont,
  keywordstyle=\bfseries,
  commentstyle=\itshape\color{codigocomentario},
  stringstyle=\color{black},
  numbers=left,
  numberstyle=\tiny\color{codigonumero},
  numbersep=8pt,
  backgroundcolor=\color{codigofundo},
  frame=tb,
  framerule=0.4pt,
  xleftmargin=1.5em,
  breaklines=true,
  breakatwhitespace=true,
  showstringspaces=false,
  tabsize=4,
  columns=fullflexible,
  keepspaces=true,
  upquote=true,
  extendedchars=true,
  inputencoding=utf8,
  % Acentos dentro de códigos e comentários (necessário com pdfLaTeX).
  literate=
    {á}{{\'a}}1 {à}{{\`a}}1 {ã}{{\~a}}1 {â}{{\^a}}1
    {é}{{\'e}}1 {ê}{{\^e}}1 {í}{{\'i}}1
    {ó}{{\'o}}1 {õ}{{\~o}}1 {ô}{{\^o}}1 {ú}{{\'u}}1 {ü}{{\"u}}1
    {ç}{{\c{c}}}1
    {Á}{{\'A}}1 {À}{{\`A}}1 {Ã}{{\~A}}1 {Â}{{\^A}}1
    {É}{{\'E}}1 {Ê}{{\^E}}1 {Í}{{\'I}}1
    {Ó}{{\'O}}1 {Õ}{{\~O}}1 {Ô}{{\^O}}1 {Ú}{{\'U}}1
    {Ç}{{\c{C}}}1
}

% ------------------------------------------------------------------------------
% Algoritmos (pacote algpseudocode) com palavras-chave em português
% ------------------------------------------------------------------------------
\algrenewcommand\algorithmicrequire{\textbf{Entrada:}}
\algrenewcommand\algorithmicensure{\textbf{Saída:}}
\algrenewcommand\algorithmicif{\textbf{se}}
\algrenewcommand\algorithmicthen{\textbf{então}}
\algrenewcommand\algorithmicelse{\textbf{senão}}
\algrenewcommand\algorithmicend{\textbf{fim}}
\algrenewcommand\algorithmicwhile{\textbf{enquanto}}
\algrenewcommand\algorithmicdo{\textbf{faça}}
\algrenewcommand\algorithmicfor{\textbf{para}}
\algrenewcommand\algorithmicforall{\textbf{para todo}}
\algrenewcommand\algorithmicrepeat{\textbf{repita}}
\algrenewcommand\algorithmicuntil{\textbf{até que}}
\algrenewcommand\algorithmicreturn{\textbf{retorne}}
\algrenewcommand\algorithmicfunction{\textbf{função}}
\algrenewcommand\algorithmicprocedure{\textbf{procedimento}}
\algrenewcommand{\algorithmiccomment}[1]{\hfill\(\triangleright\) #1}

% ------------------------------------------------------------------------------
% Gráficos (pgfplots) com vírgula decimal e ponto como separador de milhar
% ------------------------------------------------------------------------------
\pgfplotsset{
  /pgf/number format/.cd, use comma, 1000 sep={.},
}

% ------------------------------------------------------------------------------
% Referências cruzadas (cleveref)
% ------------------------------------------------------------------------------
% Use \cref{rotulo} no meio da frase e \Cref{rotulo} no início. Nos dois casos
% o nome sai com inicial maiúscula ("Figura 3", "Tabela 2"), como pede a norma.
\crefname{chapter}{\termocap}{\termocaps}
\Crefname{chapter}{\termocap}{\termocaps}
\crefname{section}{Seção}{Seções}
\Crefname{section}{Seção}{Seções}
\crefname{subsection}{Seção}{Seções}
\Crefname{subsection}{Seção}{Seções}
\crefname{subsubsection}{Seção}{Seções}
\Crefname{subsubsection}{Seção}{Seções}
\crefname{appendix}{Apêndice}{Apêndices}
\Crefname{appendix}{Apêndice}{Apêndices}
% O abnTeX2 não aciona a detecção de apêndices do cleveref. Dentro dos
% ambientes de apêndices e anexos, os capítulos passam a ser referenciados
% como "Apêndice A" e "Anexo A".
\AtBeginEnvironment{apendicesenv}{\crefalias{chapter}{appendix}}
\AtBeginEnvironment{anexosenv}{\crefalias{chapter}{anexo}}
\crefname{figure}{Figura}{Figuras}
\Crefname{figure}{Figura}{Figuras}
\crefname{subfigure}{Figura}{Figuras}
\Crefname{subfigure}{Figura}{Figuras}
\crefname{table}{Tabela}{Tabelas}
\Crefname{table}{Tabela}{Tabelas}
\crefname{quadro}{Quadro}{Quadros}
\Crefname{quadro}{Quadro}{Quadros}
\crefname{codigo}{Código}{Códigos}
\Crefname{codigo}{Código}{Códigos}
\crefname{algoritmo}{Algoritmo}{Algoritmos}
\Crefname{algoritmo}{Algoritmo}{Algoritmos}
\crefname{equation}{Equação}{Equações}
\Crefname{equation}{Equação}{Equações}
\crefname{page}{página}{páginas}
\Crefname{page}{Página}{Páginas}
\creflabelformat{equation}{(#2#1#3)}
\crefrangelabelformat{equation}{(#3#1#4) a (#5#2#6)}
\newcommand{\crefpairconjunction}{ e }
\newcommand{\crefmiddleconjunction}{, }
\newcommand{\creflastconjunction}{ e }
\newcommand{\crefrangeconjunction}{ a }

% ------------------------------------------------------------------------------
% Siglas (pacote acro)
% ------------------------------------------------------------------------------
% A lista mostra só as siglas usadas no texto, em ordem alfabética.
\acsetup{
  list/template = longtable,
  list/heading  = none,
  list/display  = used,
  list/sort     = true,
}

% ------------------------------------------------------------------------------
% Links do PDF
% ------------------------------------------------------------------------------
\definecolor{azulescuro}{RGB}{0,51,128}
\hypersetup{
  colorlinks = true,
  linkcolor  = azulescuro,
  citecolor  = azulescuro,
  urlcolor   = azulescuro,
  filecolor  = azulescuro,
  bookmarksdepth = 4,
}
% Para imprimir, links pretos ficam melhores. Descomente a linha abaixo:
% \hypersetup{hidelinks}

% ------------------------------------------------------------------------------
% Notas de revisão (aparecem só com \mostrarnotastrue no arquivo principal)
% ------------------------------------------------------------------------------
% Uso: \notaaluno{texto} e \notaorientador{texto}.
\newcommand{\notaaluno}[1]{\todo[author=\textbf{Aluno},color=green!25,inline,caption={}]{#1}}
\newcommand{\notaorientador}[1]{\todo[author=\textbf{Orientador},color=red!20,inline,caption={}]{#1}}

% ------------------------------------------------------------------------------
% Caixas de orientação e de exemplos de LaTeX
% ------------------------------------------------------------------------------
% As caixas "Orientação" explicam o que escrever em cada parte do trabalho.
% Elas somem do PDF com \mostrarinstrucoesfalse no arquivo principal.
% IMPORTANTE: \begin{instrucao} e \end{instrucao} devem ficar sozinhos na linha.
\definecolor{orientacaoborda}{RGB}{30,80,140}
\definecolor{orientacaofundo}{RGB}{240,245,252}
\definecolor{exemploborda}{gray}{0.45}

\ifmostrarinstrucoes
  \newtcolorbox{instrucao}{
    enhanced, breakable,
    colback=orientacaofundo, colframe=orientacaoborda,
    boxrule=0.6pt, arc=2pt,
    left=6pt, right=6pt, top=4pt, bottom=4pt,
    title={Orientação},
    fonttitle=\sffamily\bfseries\footnotesize,
    fontupper=\small\linespread{1}\selectfont,
    before upper={\setlength{\parindent}{0pt}\setlength{\parskip}{0.5em}},
    before skip=0.8em, after skip=0.8em,
  }
\else
  \excludecomment{instrucao}
\fi

% Subtítulo dentro de uma caixa de orientação.
\newcommand{\orientacao}[1]{{\sffamily\bfseries #1}\par}

% Código LaTeX e o resultado logo abaixo.
\newtcblisting{exemplo}{
  enhanced, listing and text,
  colback=white, colframe=exemploborda, boxrule=0.4pt, arc=1pt,
  left=4pt, right=4pt, top=2pt, bottom=2pt,
  segmentation style={dashed, exemploborda},
  title={Código LaTeX e resultado}, coltitle=black, colbacktitle=black!8,
  fonttitle=\sffamily\footnotesize,
  listing options={language={[LaTeX]TeX}, numbers=none, frame=none,
                   backgroundcolor=\color{white}, xleftmargin=0pt},
  fontlower=\small\linespread{1}\selectfont,
}

% Apenas o código LaTeX (usado para figuras, tabelas e outros objetos
% flutuantes, cujo resultado aparece no próprio texto do exemplo).
\newtcblisting{exemplocodigo}{
  enhanced, listing only,
  colback=white, colframe=exemploborda, boxrule=0.4pt, arc=1pt,
  left=4pt, right=4pt, top=2pt, bottom=2pt,
  title={Código LaTeX}, coltitle=black, colbacktitle=black!8,
  fonttitle=\sffamily\footnotesize,
  listing options={language={[LaTeX]TeX}, numbers=none, frame=none,
                   backgroundcolor=\color{white}, xleftmargin=0pt},
}

% ==============================================================================
% config/elementos.tex
% Comandos de dados e leiaute dos elementos pré-textuais: capa, folha de rosto,
% folha de aprovação, resumos e epígrafe. Os dados em si ficam em dados.tex.
% ==============================================================================

% ------------------------------------------------------------------------------
% Instituição
% ------------------------------------------------------------------------------
\newcommand{\imprimiruniversidade}{Universidade Federal do Espírito Santo}
\newcommand{\imprimircentro}{Centro Universitário Norte do Espírito Santo}
\newcommand{\imprimirdepartamento}{Departamento de Computação e Eletrônica}

% ------------------------------------------------------------------------------
% Curso: \curso{cc} (Ciência da Computação) ou \curso{ec} (Engenharia de Computação)
% ------------------------------------------------------------------------------
% Em Engenharia de Computação, o TCC se chama Projeto de Graduação (nome da
% disciplina nas normas do curso). O comando \curso ajusta o curso, o grau e
% o nome do trabalho na folha de rosto e na folha de aprovação.
\newcommand{\imprimircurso}{Ciência da Computação}
\newcommand{\imprimirgrau}{Bacharel em Ciência da Computação}
\newcommand{\imprimirtipotcc}{Trabalho de Conclusão de Curso}
\newcommand{\curso}[1]{%
  \ifthenelse{\equal{#1}{ec}}{%
    \renewcommand{\imprimircurso}{Engenharia de Computação}%
    \renewcommand{\imprimirgrau}{Bacharel em Engenharia de Computação}%
    \renewcommand{\imprimirtipotcc}{Projeto de Graduação}%
  }{%
    \renewcommand{\imprimircurso}{Ciência da Computação}%
    \renewcommand{\imprimirgrau}{Bacharel em Ciência da Computação}%
    \renewcommand{\imprimirtipotcc}{Trabalho de Conclusão de Curso}%
  }%
}
\newcommand{\imprimircolegiado}{Colegiado do Curso de \imprimircurso}

% Natureza do trabalho (folha de rosto e folha de aprovação).
\newcommand{\imprimirnatureza}{%
  \imprimirtipotcc{} apresentado ao \imprimircolegiado{} do
  \imprimircentro{} da \imprimiruniversidade, como requisito parcial para
  obtenção do grau de \imprimirgrau.}
\preambulo{\imprimirnatureza}
\tipotrabalho{\imprimirtipotcc{} (Graduação)}

\local{São Mateus, ES}

% ------------------------------------------------------------------------------
% Dados complementares
% ------------------------------------------------------------------------------
% Subtítulo (opcional). Na capa e na folha de rosto aparece após dois-pontos.
\newcommand{\imprimirsubtitulo}{}
\newcommand{\subtitulo}[1]{\renewcommand{\imprimirsubtitulo}{#1}}
\newcommand{\imprimirtitulocompleto}{%
  \imprimirtitulo\ifdefempty{\imprimirsubtitulo}{}{: \imprimirsubtitulo}}

% Data de aprovação. Deixe vazio para gerar linhas em branco a preencher.
\newcommand{\imprimirdataaprovacao}{}
\newcommand{\dataaprovacao}[1]{\renewcommand{\imprimirdataaprovacao}{#1}}

% Banca examinadora: \membrobanca{Nome com titulação}{Instituição}{Função}
\makeatletter
\newcommand{\imprimirbanca}{}
\newcommand{\membrobanca}[3]{%
  \g@addto@macro\imprimirbanca{%
    \assinatura{\textbf{#1}\\#2\ifthenelse{\equal{#3}{}}{}{\\#3}}}%
}
\makeatother

% Palavras-chave (NBR 6028:2021): separadas por ponto e vírgula e
% finalizadas por ponto. Ex.: \palavraschave{otimização; heurística; mochila}
\newcommand{\palavraschave}[1]{%
  \par\vspace{\baselineskip}\noindent\textbf{Palavras-chave}: #1.}
\newcommand{\keywords}[1]{%
  \par\vspace{\baselineskip}\noindent\textbf{Keywords}: #1.}

% ------------------------------------------------------------------------------
% Capa (NBR 14724:2024, 4.1.1)
% ------------------------------------------------------------------------------
\renewcommand{\imprimircapa}{%
  \begin{capa}%
    \centering\SingleSpacing
    \includegraphics[width=3.3cm]{figuras/logo-ufes}\par
    \vspace{0.6cm}
    {\bfseries
      \MakeTextUppercase{\imprimircentro}\par
      \MakeTextUppercase{\imprimirdepartamento}\par
      \MakeTextUppercase{\imprimircolegiado}\par}
    \vfill
    {\large\imprimirautor\par}
    \vfill
    {\bfseries\Large\imprimirtitulocompleto\par}
    \vfill\vfill
    {\large\imprimirlocal\par
     \imprimirdata\par}
  \end{capa}%
}

% ------------------------------------------------------------------------------
% Folha de rosto (NBR 14724:2024, 4.2.1.1)
% ------------------------------------------------------------------------------
\renewcommand{\imprimirfolhaderosto}{%
  \begin{folhaderosto}%
    \centering\SingleSpacing
    {\large\imprimirautor\par}
    \vfill
    {\bfseries\Large\imprimirtitulocompleto\par}
    \vfill
    \noindent\hspace*{.45\textwidth}%
    \begin{minipage}{.55\textwidth}
      \justifying\imprimirnatureza

      \vspace{\baselineskip}
      \imprimirorientadorRotulo~\imprimirorientador
      \ifdefempty{\imprimircoorientador}{}{\par
        \imprimircoorientadorRotulo~\imprimircoorientador}
    \end{minipage}\par
    \vfill\vfill
    {\large\imprimirlocal\par
     \imprimirdata\par}
  \end{folhaderosto}%
}

% ------------------------------------------------------------------------------
% Folha de aprovação (NBR 14724:2024, 4.2.1.3)
% ------------------------------------------------------------------------------
\newcommand{\imprimirfolhadeaprovacao}{%
  \begin{folhadeaprovacao}%
    \centering
    {\large\imprimirautor\par}
    \vspace{1.2cm}
    {\bfseries\Large\imprimirtitulocompleto\par}
    \vspace{1.2cm}
    \noindent\hspace*{.45\textwidth}%
    \begin{minipage}{.55\textwidth}
      \justifying\imprimirnatureza
    \end{minipage}\par
    \vspace{1cm}
    \raggedright
    Aprovado em
    \ifdefempty{\imprimirdataaprovacao}%
      {\rule{1.2cm}{0.4pt} de \rule{3.5cm}{0.4pt} de \rule{1.6cm}{0.4pt}.}%
      {\imprimirdataaprovacao.}\par
    \vspace{0.8cm}
    \centering
    {\bfseries BANCA EXAMINADORA\par}
    \imprimirbanca
  \end{folhadeaprovacao}%
}

% \justifying para minipages (o memoir não define este comando).
\providecommand{\justifying}{%
  \rightskip=0pt plus 0pt\leftskip=0pt\parfillskip=0pt plus 1fil\relax}

% ==============================================================================
% Dados do trabalho
% Preencha os campos abaixo. Eles alimentam a capa, a folha de rosto, a folha
% de aprovação e as informações do PDF.
% ==============================================================================

% Curso: cc (Ciência da Computação) ou ec (Engenharia de Computação).
% Este exemplo usa ec apenas para mostrar a opção.
\curso{ec}

% Título e subtítulo (deixe \subtitulo{} vazio se não houver).
\titulo{Monitora}
\subtitulo{um sistema web para agendamento de monitorias acadêmicas}

\autor{Nome Completo do Autor}
\data{2026}

% Orientador e coorientador, com titulação. Para orientadora, use
% \orientador[Orientadora:]{Profa. Dra. Nome}. Sem coorientador, apague a linha.
\orientador{Prof. Dr. Nome do Orientador}
% \coorientador{Prof. Me. Nome do Coorientador}

% Data da defesa. Se ficar vazio, a folha de aprovação sai com linhas em branco.
\dataaprovacao{}

% Banca examinadora: \membrobanca{Nome com titulação}{Instituição}{Função}.
% O orientador é o primeiro membro.
\membrobanca{Prof. Dr. Nome do Orientador}{Universidade Federal do Espírito Santo}{Orientador}
\membrobanca{Prof. Dr. Nome do Examinador}{Universidade Federal do Espírito Santo}{Examinador interno}
\membrobanca{Profa. Dra. Nome da Examinadora}{Instituição de origem}{Examinadora externa}

% Informações do arquivo PDF.
\makeatletter
\hypersetup{
  pdftitle   = {\@title},
  pdfauthor  = {\@author},
  pdfsubject = {\imprimirtipotrabalho},
  pdfkeywords= {sistema web; agendamento; monitoria; engenharia de software},
}
\makeatother

% ==============================================================================
% Siglas e abreviaturas (pacote acro)
% ------------------------------------------------------------------------------
% Declare aqui cada sigla uma única vez. No texto, use:
%   \ac{chave}   primeira ocorrência: forma extensa seguida da sigla entre
%                parênteses; nas seguintes, só a sigla;
%   \acs{chave}  sempre a sigla (use em títulos e legendas);
%   \acl{chave}  sempre a forma extensa;
%   \acp{chave}  plural da sigla.
% A lista de abreviaturas e siglas é montada automaticamente, em ordem
% alfabética, apenas com as siglas usadas no texto.
%
% Campos:
%   short   sigla;
%   long    forma extensa usada no texto;
%   list    forma extensa na lista (opcional, para ajustar maiúsculas);
%   foreign forma original, quando a sigla vem de outro idioma.
% ==============================================================================

\DeclareAcronym{rf}{
  short = RF,
  long  = requisito funcional,
  list  = Requisito funcional,
  long-plural-form = requisitos funcionais,
}

\DeclareAcronym{rnf}{
  short = RNF,
  long  = requisito não funcional,
  list  = Requisito não funcional,
  long-plural-form = requisitos não funcionais,
}

\DeclareAcronym{mtv}{
  short   = MTV,
  long    = {modelo, \emph{template} e visão},
  list    = {Modelo, \emph{template} e visão},
  foreign = Model-Template-View,
  foreign-babel = english,
}

\DeclareAcronym{orm}{
  short   = ORM,
  long    = mapeamento objeto-relacional,
  list    = Mapeamento objeto-relacional,
  foreign = Object-Relational Mapping,
  foreign-babel = english,
}

\DeclareAcronym{sus}{
  short   = SUS,
  long    = escala de usabilidade de sistemas,
  list    = Escala de usabilidade de sistemas,
  foreign = System Usability Scale,
  foreign-babel = english,
}

\DeclareAcronym{uml}{
  short   = UML,
  long    = linguagem de modelagem unificada,
  list    = Linguagem de modelagem unificada,
  foreign = Unified Modeling Language,
  foreign-babel = english,
}

\DeclareAcronym{ia}{
  short = IA,
  long  = inteligência artificial,
  list  = Inteligência artificial,
}

\DeclareAcronym{abnt}{
  short = ABNT,
  long  = Associação Brasileira de Normas Técnicas,
}

\addbibresource{referencias.bib}

% ==============================================================================
\begin{document}
% ==============================================================================

% ------------------------------------------------------------------------------
% Elementos pré-textuais
% ------------------------------------------------------------------------------
% A capa não é contada. A contagem começa na folha de rosto (página 1), mas os
% números só aparecem a partir da Introdução.
\pagenumbering{Alph}
\imprimircapa
\pagenumbering{arabic}

\imprimirfolhaderosto
\imprimirfolhadeaprovacao

% Elementos opcionais: comente as linhas que não for usar.
% ==============================================================================
% Dedicatória (opcional, NBR 14724:2024, 4.2.1.4)
% Para não usar, comente a linha % ==============================================================================
% Dedicatória (opcional, NBR 14724:2024, 4.2.1.4)
% Para não usar, comente a linha % ==============================================================================
% Dedicatória (opcional, NBR 14724:2024, 4.2.1.4)
% Para não usar, comente a linha \input{comum/dedicatoria} no arquivo principal.
% ==============================================================================
\begin{dedicatoria}
\begin{instrucao}
\orientacao{Dedicatória}
Elemento opcional, sem título, em que você presta homenagem ou dedica o
trabalho a alguém. Costuma ter uma ou duas frases, posicionadas na parte
inferior da página.
\end{instrucao}
\vspace*{\fill}
\begin{flushright}
  \begin{minipage}{.5\textwidth}
    \raggedleft\itshape
    Aos meus pais, que sempre apoiaram meus estudos.
  \end{minipage}
\end{flushright}
\end{dedicatoria}
 no arquivo principal.
% ==============================================================================
\begin{dedicatoria}
\begin{instrucao}
\orientacao{Dedicatória}
Elemento opcional, sem título, em que você presta homenagem ou dedica o
trabalho a alguém. Costuma ter uma ou duas frases, posicionadas na parte
inferior da página.
\end{instrucao}
\vspace*{\fill}
\begin{flushright}
  \begin{minipage}{.5\textwidth}
    \raggedleft\itshape
    Aos meus pais, que sempre apoiaram meus estudos.
  \end{minipage}
\end{flushright}
\end{dedicatoria}
 no arquivo principal.
% ==============================================================================
\begin{dedicatoria}
\begin{instrucao}
\orientacao{Dedicatória}
Elemento opcional, sem título, em que você presta homenagem ou dedica o
trabalho a alguém. Costuma ter uma ou duas frases, posicionadas na parte
inferior da página.
\end{instrucao}
\vspace*{\fill}
\begin{flushright}
  \begin{minipage}{.5\textwidth}
    \raggedleft\itshape
    Aos meus pais, que sempre apoiaram meus estudos.
  \end{minipage}
\end{flushright}
\end{dedicatoria}

% ==============================================================================
% Agradecimentos (opcional, NBR 14724:2024, 4.2.1.5)
% Para não usar, comente a linha % ==============================================================================
% Agradecimentos (opcional, NBR 14724:2024, 4.2.1.5)
% Para não usar, comente a linha % ==============================================================================
% Agradecimentos (opcional, NBR 14724:2024, 4.2.1.5)
% Para não usar, comente a linha \input{comum/agradecimentos} no arquivo principal.
% ==============================================================================
\begin{agradecimentos}
\begin{instrucao}
\orientacao{Agradecimentos}
Elemento opcional, dirigido a quem contribuiu de forma relevante para o
trabalho. Seja breve (meia página costuma bastar). Se o trabalho recebeu bolsa
ou auxílio, informe aqui a agência de fomento e o número do processo, porque
muitas agências exigem esse registro.
\end{instrucao}

Agradeço ao meu orientador pela paciência e pelas discussões que deram forma a
este trabalho.

Aos professores do Departamento de Computação e Eletrônica, pela formação que
recebi ao longo do curso.

À minha família e aos colegas de turma, pelo apoio nos momentos mais difíceis.
\end{agradecimentos}
 no arquivo principal.
% ==============================================================================
\begin{agradecimentos}
\begin{instrucao}
\orientacao{Agradecimentos}
Elemento opcional, dirigido a quem contribuiu de forma relevante para o
trabalho. Seja breve (meia página costuma bastar). Se o trabalho recebeu bolsa
ou auxílio, informe aqui a agência de fomento e o número do processo, porque
muitas agências exigem esse registro.
\end{instrucao}

Agradeço ao meu orientador pela paciência e pelas discussões que deram forma a
este trabalho.

Aos professores do Departamento de Computação e Eletrônica, pela formação que
recebi ao longo do curso.

À minha família e aos colegas de turma, pelo apoio nos momentos mais difíceis.
\end{agradecimentos}
 no arquivo principal.
% ==============================================================================
\begin{agradecimentos}
\begin{instrucao}
\orientacao{Agradecimentos}
Elemento opcional, dirigido a quem contribuiu de forma relevante para o
trabalho. Seja breve (meia página costuma bastar). Se o trabalho recebeu bolsa
ou auxílio, informe aqui a agência de fomento e o número do processo, porque
muitas agências exigem esse registro.
\end{instrucao}

Agradeço ao meu orientador pela paciência e pelas discussões que deram forma a
este trabalho.

Aos professores do Departamento de Computação e Eletrônica, pela formação que
recebi ao longo do curso.

À minha família e aos colegas de turma, pelo apoio nos momentos mais difíceis.
\end{agradecimentos}

% ==============================================================================
% Epígrafe (opcional, NBR 14724:2024, 4.2.1.6)
% Para não usar, comente a linha % ==============================================================================
% Epígrafe (opcional, NBR 14724:2024, 4.2.1.6)
% Para não usar, comente a linha % ==============================================================================
% Epígrafe (opcional, NBR 14724:2024, 4.2.1.6)
% Para não usar, comente a linha \input{comum/epigrafe} no arquivo principal.
% ==============================================================================
\begin{epigrafe}
\begin{instrucao}
\orientacao{Epígrafe}
Elemento opcional, sem título, com uma citação relacionada ao tema do
trabalho. A NBR 14724:2024 fixa a posição na parte inferior da página, ocupando
da metade da mancha gráfica até a margem direita. A citação segue a NBR
10520:2023 (autoria, ano e página) e a obra citada deve constar nas
referências. Textos traduzidos por você levam a expressão ``tradução nossa''.
\end{instrucao}
\vspace*{\fill}
\hfill
\begin{minipage}{.5\textwidth}
  \SingleSpacing
  ``A Máquina Analítica não tem a pretensão de criar coisa alguma. Ela pode
  fazer qualquer coisa que saibamos ordenar que ela execute''
  \cite[p.~722, tradução nossa]{lovelace1843sketch}.
\end{minipage}
\end{epigrafe}
 no arquivo principal.
% ==============================================================================
\begin{epigrafe}
\begin{instrucao}
\orientacao{Epígrafe}
Elemento opcional, sem título, com uma citação relacionada ao tema do
trabalho. A NBR 14724:2024 fixa a posição na parte inferior da página, ocupando
da metade da mancha gráfica até a margem direita. A citação segue a NBR
10520:2023 (autoria, ano e página) e a obra citada deve constar nas
referências. Textos traduzidos por você levam a expressão ``tradução nossa''.
\end{instrucao}
\vspace*{\fill}
\hfill
\begin{minipage}{.5\textwidth}
  \SingleSpacing
  ``A Máquina Analítica não tem a pretensão de criar coisa alguma. Ela pode
  fazer qualquer coisa que saibamos ordenar que ela execute''
  \cite[p.~722, tradução nossa]{lovelace1843sketch}.
\end{minipage}
\end{epigrafe}
 no arquivo principal.
% ==============================================================================
\begin{epigrafe}
\begin{instrucao}
\orientacao{Epígrafe}
Elemento opcional, sem título, com uma citação relacionada ao tema do
trabalho. A NBR 14724:2024 fixa a posição na parte inferior da página, ocupando
da metade da mancha gráfica até a margem direita. A citação segue a NBR
10520:2023 (autoria, ano e página) e a obra citada deve constar nas
referências. Textos traduzidos por você levam a expressão ``tradução nossa''.
\end{instrucao}
\vspace*{\fill}
\hfill
\begin{minipage}{.5\textwidth}
  \SingleSpacing
  ``A Máquina Analítica não tem a pretensão de criar coisa alguma. Ela pode
  fazer qualquer coisa que saibamos ordenar que ela execute''
  \cite[p.~722, tradução nossa]{lovelace1843sketch}.
\end{minipage}
\end{epigrafe}


% Resumos (obrigatórios).
% ==============================================================================
% Resumo na língua vernácula (obrigatório, NBR 14724:2024 e NBR 6028:2021)
% ==============================================================================
\begin{resumo}
\begin{instrucao}
\orientacao{Resumo}
O resumo apresenta, de forma concisa, o objetivo, o método, os resultados e as
conclusões do trabalho. Escreva-o por último, quando o texto estiver pronto.
Siga a NBR 6028:2021:
\begin{itemize}
  \item parágrafo único, com 150 a 500 palavras;
  \item frases afirmativas, com o verbo na terceira pessoa, como recomenda o
        item 4.1.4 da norma (``Este trabalho propõe\ldots'', ``Os resultados
        mostram\ldots''). A voz ativa deixou de ser exigida, mas costuma
        deixar o resumo mais direto;
  \item sem citações, fórmulas, siglas não explicadas, figuras ou tabelas;
  \item palavras-chave logo abaixo, separadas por ponto e vírgula e finalizadas
        por ponto. Prefira termos que um leitor usaria para encontrar o
        trabalho num catálogo (de três a cinco é a prática mais comum).
\end{itemize}
Um roteiro que funciona bem dedica uma ou duas frases ao contexto e ao
problema, outras ao objetivo e ao método e as últimas aos principais
resultados (com números) e à conclusão.
\end{instrucao}

A monitoria acadêmica é uma das principais formas de apoio aos estudantes nas
disciplinas iniciais dos cursos de Computação, mas a marcação de atendimentos
costuma depender de mensagens informais e planilhas, o que gera conflitos de
horário e dificulta o acompanhamento pela coordenação. Este trabalho apresenta
o Monitora, um sistema web para agendamento de monitorias que permite aos
monitores publicar horários de atendimento, aos estudantes reservar esses
horários e aos professores acompanhar a frequência e os temas atendidos. O
sistema foi desenvolvido com um processo iterativo baseado no Scrum, em seis
ciclos de duas semanas. Os requisitos foram levantados em entrevistas com
monitores e professores e especificados por meio de histórias de usuário. A
implementação usa o arcabouço Django, o banco de dados PostgreSQL e
contêineres Docker. A verificação combinou testes automatizados, com cobertura
de 87\% do código, e uma avaliação com 18 usuários, que atribuíram ao sistema
pontuação média de 84,6 na escala de usabilidade de sistemas, próxima da faixa
considerada excelente. Os resultados indicam que o sistema atende aos requisitos
levantados e que pode substituir o processo informal usado atualmente nas
disciplinas do curso.

\palavraschave{sistema web; agendamento; monitoria acadêmica; engenharia de
software}
\end{resumo}

% ==============================================================================
% Resumo em língua estrangeira (obrigatório, NBR 14724:2024)
% ==============================================================================
\begin{resumo}[Abstract]
\begin{instrucao}
\orientacao{Abstract}
Versão do resumo em inglês, com o mesmo conteúdo e as mesmas regras. Não use
tradução automática sem revisão cuidadosa. Termos técnicos devem seguir o uso
da área (por exemplo, \emph{user story} para ``história de usuário'', e não uma
tradução literal). Se usar ferramenta de IA para traduzir ou revisar, declare
o uso no apêndice sobre inteligência artificial.
\end{instrucao}

\begin{otherlanguage*}{english}
Academic tutoring is one of the main forms of student support in the initial
courses of Computing programs, but scheduling sessions usually relies on
informal messages and spreadsheets, which leads to schedule conflicts and
makes it difficult for coordinators to follow up. This work presents Monitora,
a web system for scheduling tutoring sessions that allows tutors to publish
their available hours, students to book those hours and professors to track
attendance and the topics covered. The system was developed with an iterative
process based on Scrum, in six two-week cycles. Requirements were gathered in
interviews with tutors and professors and specified as user stories. The
implementation uses the Django framework, the PostgreSQL database and Docker
containers. Verification combined automated tests, with 87\% code coverage,
and an evaluation with 18 users, who gave the system an average score of 84.6
on the System Usability Scale, close to the range rated as excellent. The results indicate that
the system meets the gathered requirements and can replace the informal
process currently used in the program's courses.

\keywords{web system; scheduling; academic tutoring; software engineering}
\end{otherlanguage*}
\end{resumo}


% Listas: mantenha apenas as dos tipos de ilustração que o trabalho usa.
% Lista vazia não deve aparecer no trabalho.
\pdfbookmark[0]{\listfigurename}{lof}
\listoffigures*
\cleardoublepage

\pdfbookmark[0]{\listofquadrosname}{loq}
\listofquadros*
\cleardoublepage

\pdfbookmark[0]{\listtablename}{lot}
\listoftables*
\cleardoublepage

\pdfbookmark[0]{\listofcodigosname}{lcod}
\listofcodigos*
\cleardoublepage

% \pdfbookmark[0]{\listofalgoritmosname}{lalg}
% \listofalgoritmos*
% \cleardoublepage

% Lista de abreviaturas e siglas (gerada a partir de tecnologico/siglas.tex).
\pretextualchapter{\listadesiglasname}
\printacronyms
\cleardoublepage

% Sumário.
\pdfbookmark[0]{\contentsname}{toc}
\tableofcontents*
\cleardoublepage

% ------------------------------------------------------------------------------
% Elementos textuais
% ------------------------------------------------------------------------------
\textual

% ==============================================================================
% Capítulo 1: Introdução
% ==============================================================================
\chapter{Introdução}
\label{cap:introducao}

\begin{instrucao}
\orientacao{O que vai na Introdução}
A Introdução situa quem lê o trabalho. Em um trabalho tecnológico, cujo foco principal é
o produto desenvolvido, ela apresenta o contexto, o problema que o software
resolve, quem são os usuários, os objetivos, a justificativa, o método de
desenvolvimento e a organização do texto. Como referência, de três a cinco
páginas costumam ser suficientes. Escreva uma primeira versão no início e
reescreva no final, porque a Introdução promete o que o resto do texto precisa
cumprir.

Comece pelo geral e avance até o problema específico. Descreva o problema do
ponto de vista de quem o enfrenta, e não da tecnologia. ``Não existe um
aplicativo em Django para monitorias'' não é um problema. ``Estudantes não
conseguem saber quando há monitor disponível'' é.

\orientacao{Capítulo ou seção?}
A NBR 14724:2024 determina que o trabalho seja dividido em seções, e não em
capítulos. No LaTeX nada muda (continue usando \verb|\chapter|, que gera o
título ``1 INTRODUÇÃO''). A diferença está só no nome usado nas referências
cruzadas, que este modelo controla pelo comando \verb|\termocap| em
\verb|config/estilo.tex|. O padrão do modelo é ``Capítulo'', o termo mais
usado nos cursos.

\orientacao{Rótulos}
Logo após cada \verb|\chapter| ou \verb|\section|, coloque um
\verb|\label| com um prefixo que indique o tipo: \verb|cap:| para capítulos,
\verb|sec:| para seções, \verb|fig:| para figuras, \verb|tab:| para tabelas,
\verb|qua:| para quadros, \verb|cod:| para códigos e assim por diante. Nunca
use acentos ou espaços nos rótulos.
\end{instrucao}

\begin{instrucao}
\orientacao{Pessoa verbal e estilo acadêmico}
Textos acadêmicos são escritos, por convenção, de forma impessoal. A NBR
14724:2024 não determina a pessoa verbal do trabalho, mas a NBR 6028:2021
recomenda o verbo na terceira pessoa no resumo (item 4.1.4). Este modelo
estende essa recomendação a todo o texto, para que ele fique uniforme. Use:
\begin{itemize}
  \item a terceira pessoa do singular, com o trabalho, o método ou o sistema
        como sujeito (``este trabalho propõe'', ``a heurística obteve'',
        ``o sistema permite'');
  \item ou a voz passiva (``os experimentos foram executados'', ``observa-se
        que'', ``adotou-se'').
\end{itemize}
Evite a primeira pessoa do singular (``eu proponho'') e do plural (``nós
propomos'', ``vamos mostrar''), mesmo como plural de modéstia. Algumas áreas
aceitam o ``nós'', mas a forma impessoal é a convenção predominante nos
trabalhos acadêmicos. Qualquer que seja a escolha, não misture as duas formas
no mesmo texto, um erro frequente nos TCCs. As exceções são
os agradecimentos, a dedicatória e a declaração de uso de inteligência
artificial, escritos na primeira pessoa do singular por serem manifestações
pessoais do autor.

\begin{itemize}
  \item \textbf{Evite:} ``Neste trabalho, vamos propor uma heurística e
        mostraremos que ela é melhor.''
  \item \textbf{Prefira:} ``Este trabalho propõe uma heurística e avalia seu
        desempenho em relação a \ldots''
\end{itemize}

Algumas regras de redação ajudam o texto a soar claro e profissional:
\begin{itemize}
  \item escreva frases diretas e evite expressões de preenchimento como ``vale
        ressaltar que'' e ``é importante notar que'';
  \item evite adjetivos de efeito (``crucial'', ``poderoso'', ``robusto'') e
        afirmações sem dados ou citação (``muito usado'', ``o melhor método'');
  \item não use linguagem coloquial nem perguntas retóricas;
  \item use vírgula decimal e ponto como separador de milhar (0,41\% e
        1.200), inclusive em tabelas e gráficos.
\end{itemize}
\end{instrucao}

A monitoria é uma das principais ações de apoio ao estudante nas disciplinas
iniciais dos cursos de Computação. Nela, estudantes que já cursaram a
disciplina atendem colegas com dificuldades, sob a supervisão de um professor.
Nas disciplinas de programação e de estrutura de dados, em que as taxas de
reprovação costumam ser altas, a monitoria é frequentemente o primeiro recurso
procurado pelos alunos.

No curso usado como estudo de caso, os horários de monitoria são divulgados em
grupos de mensagens e a marcação de atendimentos depende de conversas
individuais com os monitores. Esse processo informal gera conflitos de
horário, atendimentos não registrados e falta de informação para o professor
supervisor, que precisa comprovar a frequência dos monitores e não sabe quais
assuntos geram mais dúvidas.

\begin{instrucao}
\orientacao{Siglas}
Declare cada sigla uma única vez no arquivo \verb|tecnologico/siglas.tex| e use
\verb|\ac{chave}| no texto. Na primeira ocorrência o LaTeX escreve a forma
extensa seguida da sigla. Nas seguintes, apenas a sigla. A lista de siglas é
montada sozinha, só com as siglas usadas.
\begin{exemplo}
A \ac{abnt} publica as normas.
Este trabalho segue as
normas da \ac{abnt}.
\end{exemplo}
\acreset{abnt}% desfaz o uso da sigla no exemplo acima
Em títulos e legendas use \verb|\acs{chave}|, que imprime só a sigla. Para o
plural, use \verb|\acp{chave}|. Não use \verb|\ac| no resumo.
\end{instrucao}

% ------------------------------------------------------------------------------
\section{Problema}
\label{sec:introducao:problema}

O problema tratado neste trabalho é a ausência de um meio centralizado para
divulgar os horários de monitoria, reservar atendimentos e registrar o que foi
atendido. A falta desse meio afeta três grupos de usuários. Os estudantes não
sabem quando há monitor disponível e perdem tempo procurando atendimento. Os
monitores recebem pedidos fora do horário e não têm como organizar a agenda. Os
professores não conseguem acompanhar a frequência dos monitores nem identificar
os conteúdos com mais dúvidas.

% ------------------------------------------------------------------------------
\section{Objetivos}
\label{sec:introducao:objetivos}

\begin{instrucao}
\orientacao{Como escrever os objetivos}
O \textbf{objetivo geral} expressa o resultado principal do trabalho em uma
única frase iniciada por verbo no infinitivo. Em um trabalho tecnológico, ele
costuma ser ``desenvolver um sistema que\ldots'', seguido do que o sistema
permite fazer e para quem. Os \textbf{objetivos específicos} são os resultados
intermediários e devem ser verificáveis. Em software, eles costumam
corresponder às etapas do desenvolvimento (levantar requisitos, projetar a
arquitetura, implementar, avaliar com usuários).

Evite verbos vagos como \emph{estudar}, \emph{conhecer} ou \emph{aprender},
que descrevem o processo de estudo, e não o resultado do trabalho. Evite também
objetivos que são apenas tecnologias (``usar o Django'').
\begin{itemize}
  \item \textbf{Fraco:} ``Aprender Django e fazer um site de monitorias.''
  \item \textbf{Melhor:} ``Desenvolver um sistema web que permita a estudantes
        reservar horários de monitoria e a professores acompanhar os
        atendimentos.''
\end{itemize}
Para listas numeradas por letras, como pede a NBR 6024 para alíneas, use o
ambiente \verb|alineas| do abnTeX2.
\end{instrucao}

O objetivo geral deste trabalho é desenvolver um sistema web que permita aos
monitores publicar seus horários de atendimento, aos estudantes reservar esses
horários e aos professores acompanhar a frequência e os temas das monitorias.
Para alcançá-lo, foram definidos os seguintes objetivos específicos:
\begin{alineas}
  \item levantar e especificar os requisitos do sistema com monitores e
        professores do curso;
  \item comparar o sistema proposto com as soluções usadas atualmente e com
        ferramentas de agendamento disponíveis no mercado;
  \item projetar a arquitetura e o modelo de dados do sistema;
  \item implementar o sistema e implantá-lo em ambiente de homologação;
  \item verificar o sistema por meio de testes automatizados e avaliar sua
        usabilidade com usuários reais.
\end{alineas}

% ------------------------------------------------------------------------------
\section{Justificativa}
\label{sec:introducao:justificativa}

\begin{instrucao}
\orientacao{Justificativa e citações}
Mostre por que o sistema é necessário e por que as soluções existentes não
bastam. Dados concretos ajudam (número de monitores, de atendimentos, de
estudantes atendidos), assim como relatos dos usuários levantados nas
entrevistas. Toda afirmação que não seja de conhecimento geral da área precisa
de citação.

Este modelo usa o biblatex-abnt, que segue a NBR 10520:2023. Na nova norma, o
sobrenome do autor aparece com apenas a inicial maiúscula também dentro dos
parênteses. Use \verb|\textcite| quando o autor faz parte da frase e
\verb|\cite| quando a citação vai entre parênteses.
\begin{exemplo}
Segundo \textcite{cohn2004user},
histórias devem ser pequenas.

Histórias devem ser pequenas
\cite{cohn2004user}.

Com página:
\cite[p.~23]{cohn2004user}.

Várias obras:
\cite{sommerville2018engenharia,
pressman2021engenharia}.
\end{exemplo}
As chaves (como \verb|cohn2004user|) são definidas no arquivo
\verb|referencias.bib|, que traz exemplos de livros, artigos, capítulos,
páginas web e ferramentas. Só entram nas Referências as obras citadas no texto.
\end{instrucao}

Ferramentas genéricas de agendamento resolvem parte do problema, mas não
registram a frequência exigida nos relatórios de monitoria nem associam os
atendimentos às disciplinas e aos conteúdos. Planilhas compartilhadas, por sua
vez, não impedem reservas simultâneas no mesmo horário e dependem de
preenchimento manual. Um sistema próprio, construído a partir das necessidades
levantadas com os usuários, permite tratar essas lacunas e gerar dados para a
gestão das monitorias pelo colegiado.

Do ponto de vista da formação, o desenvolvimento do sistema integra conteúdos
de engenharia de requisitos, projeto de software, bancos de dados e testes,
aplicados a um problema real da comunidade acadêmica
\cite{sommerville2018engenharia}.

% ------------------------------------------------------------------------------
\section{Método de desenvolvimento}
\label{sec:introducao:metodo}

\begin{instrucao}
\orientacao{Método de desenvolvimento}
Descreva o processo seguido (Scrum, Kanban, processo incremental), como os
requisitos foram levantados, como o trabalho foi dividido em ciclos e como o
progresso foi acompanhado. Se fez estudo de viabilidade, cronograma e análise
de riscos (previstos nas competências do TCC de Ciência da Computação e úteis
em qualquer projeto), resuma-os aqui e coloque os detalhes em um apêndice,
como no \cref{ap:planejamento}.
\end{instrucao}

O sistema foi desenvolvido com um processo iterativo baseado no Scrum
\cite{schwaber2020scrum}, adaptado a uma equipe de uma pessoa. O trabalho foi
dividido em seis ciclos (\emph{sprints}) de duas semanas, cada um com
planejamento, implementação, revisão com o orientador e retrospectiva. Os
requisitos foram levantados em entrevistas com quatro monitores e dois
professores supervisores e registrados como histórias de usuário. O
cronograma e a análise de riscos do projeto estão no \cref{ap:planejamento}.

% ------------------------------------------------------------------------------
\section{Organização do trabalho}
\label{sec:introducao:organizacao}

\begin{instrucao}
\orientacao{Organização do texto}
Feche a Introdução descrevendo o que cada capítulo apresenta. Use
\verb|\Cref{rotulo}| no início da frase e \verb|\cref{rotulo}| no meio dela. O
número é atualizado sozinho se a ordem dos capítulos mudar.
\end{instrucao}

Além desta introdução, o trabalho está organizado da seguinte forma. O
\cref{cap:fundamentacao} apresenta os conceitos e as tecnologias usados no
desenvolvimento. O \cref{cap:solucoes} compara o sistema proposto com
soluções existentes. O \cref{cap:requisitos} apresenta os requisitos do
sistema. O \cref{cap:projeto} descreve a arquitetura e os modelos de
software. O \cref{cap:implementacao} trata da implementação e o
\cref{cap:testes} apresenta a verificação e a avaliação com usuários. Por fim, o \cref{cap:conclusao}
apresenta as considerações finais.

\begin{instrucao}
\orientacao{Projetos de hardware e projetos técnicos profissionais}
Este modelo traz o exemplo de um projeto de software. As normas de Engenharia
de Computação admitem também projetos de hardware e projetos técnicos
profissionais. Nesses casos, aproveite a estrutura do modelo e adapte os
capítulos que vão dos requisitos aos testes. Para um projeto de hardware, uma
organização possível é:
\begin{enumerate}
  \item Introdução;
  \item Fundamentação teórica e tecnologias (princípios físicos, sensores,
        microcontroladores, protocolos de comunicação);
  \item Soluções relacionadas (produtos e protótipos semelhantes, com quadro
        comparativo de funcionalidades e custos);
  \item Especificação de requisitos (funcionais, elétricos, de consumo de
        energia, ambientais e de custo);
  \item Projeto do hardware (diagrama de blocos, seleção de componentes,
        esquemáticos e arquitetura do \emph{firmware});
  \item Implementação (montagem, placa de circuito impresso, programação do
        \emph{firmware} e fotos do protótipo);
  \item Testes e medições (procedimentos, instrumentos usados, resultados e
        comparação com os requisitos);
  \item Considerações finais.
\end{enumerate}
Inclua uma estimativa de custos, com um quadro de componentes, quantidades e
preços (informe a data da cotação). Esquemáticos completos e listas de
materiais vão para os apêndices. Exporte os esquemáticos do KiCad, ou da
ferramenta que usar, em PDF, para manter a qualidade vetorial.

Um projeto técnico profissional segue a lógica de um projeto de engenharia.
Ele parte do problema e dos requisitos, compara alternativas e apresenta a
solução escolhida com memorial descritivo e de cálculo, especificações, custos
e cronograma de execução. Defina a estrutura com o orientador, a partir dos
documentos que esse tipo de projeto exige na prática profissional.
\end{instrucao}

% ==============================================================================
% Capítulo 2: Fundamentação teórica e tecnologias
% ==============================================================================
\chapter{Fundamentação teórica e tecnologias}
\label{cap:fundamentacao}

\begin{instrucao}
\orientacao{O que vai na Fundamentação}
Em um trabalho tecnológico, este capítulo reúne dois tipos de conteúdo:
\begin{itemize}
  \item os \textbf{conceitos} necessários para entender o desenvolvimento
        (processos, técnicas de requisitos, padrões de arquitetura, critérios
        de qualidade e de usabilidade), com citações de livros e artigos;
  \item as \textbf{tecnologias} usadas (linguagens, arcabouços, bancos de
        dados, serviços), com uma descrição curta e a justificativa da escolha.
\end{itemize}
Não transforme o capítulo em tutorial. A banca não precisa saber instalar o
Django, mas precisa entender por que ele foi escolhido e que papel cumpre no
sistema. Inclua apenas o que será usado nos capítulos seguintes. Como
referência, de oito a doze páginas costumam ser suficientes.

Escreva com suas palavras e cite as fontes. Copiar texto sem citar é plágio e,
pelas regras do TCC, leva à reprovação com nota zero. Documentação oficial de
tecnologias pode ser citada como página web (veja a entrada
\verb|django2026documentation| em \verb|referencias.bib|).
\end{instrucao}

Este capítulo apresenta os conceitos e as tecnologias que fundamentam o
desenvolvimento do sistema. A \cref{sec:fundamentacao:requisitos} trata da
especificação de requisitos com histórias de usuário. A
\cref{sec:fundamentacao:arquitetura} apresenta o padrão de arquitetura
adotado, a \cref{sec:fundamentacao:usabilidade} discute a avaliação de
usabilidade e a \cref{sec:fundamentacao:tecnologias} descreve as tecnologias
utilizadas.

% ------------------------------------------------------------------------------
\section{Requisitos e histórias de usuário}
\label{sec:fundamentacao:requisitos}

Requisitos descrevem o que um sistema deve fazer e as restrições sob as quais
deve operar \cite{sommerville2018engenharia}. Eles costumam ser divididos em
\acp{rf}, que descrevem serviços e comportamentos do sistema, e \acp{rnf}, que
tratam de propriedades como desempenho, segurança, usabilidade e
manutenibilidade \cite{sommerville2018engenharia}.

Em processos ágeis, os requisitos funcionais são frequentemente registrados
como histórias de usuário, frases curtas escritas do ponto de vista de quem
usará o sistema, no formato ``como \emph{papel}, quero \emph{ação} para
\emph{benefício}'' \cite{cohn2004user}. Cada história é acompanhada de
critérios de aceitação, que definem quando ela está concluída.
\textcite{cohn2004user} recomenda que boas histórias sejam independentes,
negociáveis, valiosas para o usuário, estimáveis, pequenas e testáveis.

\begin{instrucao}
\orientacao{Citação direta}
A citação direta reproduz literalmente o texto da fonte e sempre informa a
página. Com até três linhas, fica no próprio parágrafo, entre aspas duplas.
\begin{exemplo}
A portaria define alucinação como
``informação aparentemente plausível,
porém incorreta, imprecisa ou
inventada'' \cite[p.~1]{ufes2026portaria}.
\end{exemplo}
Com mais de três linhas, use o ambiente \verb|citacao|, que aplica recuo de
4~cm, fonte menor e espaçamento simples, sem aspas.
\begin{exemplo}
\begin{citacao}
O usuário é responsável pela precisão,
pela originalidade e pela adequação dos
conteúdos gerados ou modificados por
sistema de inteligência artificial que
utilizar, devendo verificá-los antes de
sua utilização, considerada a
possibilidade de alucinação
\cite[p.~5]{ufes2026portaria}.
\end{citacao}
\end{exemplo}
Use citações diretas com moderação. Em uma Fundamentação bem escrita, a maior
parte das citações é indireta (paráfrase com indicação da fonte).
\end{instrucao}

% ------------------------------------------------------------------------------
\section{Arquitetura de software e o padrão MTV}
\label{sec:fundamentacao:arquitetura}

A arquitetura de um sistema é o conjunto de estruturas necessárias para
raciocinar sobre ele, formado pelos elementos de software, pelas relações
entre eles e pelas propriedades de ambos \cite{bass2021software}. Em
aplicações web, um padrão recorrente é separar a lógica de dados, a lógica de
apresentação e o tratamento das requisições.

O arcabouço Django organiza as aplicações segundo o padrão \ac{mtv}
\cite{django2026documentation}. O modelo define a estrutura dos dados e as
regras de negócio associadas a eles e é mapeado para tabelas do banco de dados
por meio de um \ac{orm}. A visão recebe a requisição, consulta os modelos e
escolhe a resposta. O \emph{template} define como os dados são apresentados
ao usuário. A \cref{fig:mtv} mostra o caminho de uma requisição nesse padrão.

\begin{figure}[htbp]
  \centering
  \caption{Fluxo de uma requisição no padrão \acs{mtv} do Django}
  \label{fig:mtv}
  \begin{tikzpicture}[
      node distance=8mm and 14mm,
      caixa/.style={draw, rounded corners=2pt, minimum width=2.6cm,
                    minimum height=9mm, align=center, font=\footnotesize},
      seta/.style={-{Stealth}, thick},
    ]
    \node[caixa] (nav) {Navegador};
    \node[caixa, right=of nav] (url) {Roteador de URLs\\(\texttt{urls.py})};
    \node[caixa, right=of url] (visao) {Visão\\(\texttt{views.py})};
    \node[caixa, below=of visao] (modelo) {Modelo\\(\texttt{models.py})};
    \node[caixa, below=of modelo, fill=gray!10] (bd) {Banco de dados};
    \node[caixa, below=of url] (tpl) {\emph{Template}\\(HTML)};
    \draw[seta] (nav) -- node[above, font=\scriptsize] {requisição} (url);
    \draw[seta] (url) -- (visao);
    \draw[seta] (visao) -- (modelo);
    \draw[seta] (modelo) -- node[right, font=\scriptsize] {\acs{orm}} (bd);
    \draw[seta] (visao) -- (tpl);
    \draw[seta] (tpl.west) -| node[below, pos=0.25, font=\scriptsize] {resposta} (nav.south);
  \end{tikzpicture}
  \fonte{Elaborado pelo próprio autor com base em \textcite{django2026documentation}.}
\end{figure}

\begin{instrucao}
\orientacao{Figuras: inclusão, citação e qualidade}
Pela NBR 14724, a identificação da figura (a palavra ``Figura'', o número e
o título) fica \textbf{acima} dela e a fonte fica \textbf{abaixo}, mesmo
quando a figura foi feita por você. Por isso, no LaTeX, o \verb|\caption| vem
antes do conteúdo e o \verb|\fonte| depois. Para material próprio, use
\verb|\fonteautor|. Para figura baseada em outra obra, escreva, por exemplo,
\verb|\fonte{Adaptado de \textcite{chave}.}|.
\begin{exemplocodigo}
\begin{figure}[htbp]
  \centering
  \caption{Tela de reserva de horários}
  \label{fig:tela-reserva}
  \includegraphics[width=0.9\textwidth]{figuras/tecnologico/tela-reserva.png}
  \fonteautor
\end{figure}
\end{exemplocodigo}

\textbf{Toda figura, tabela, quadro ou código incluído no trabalho deve ser
citado no texto pelo menos uma vez}, de preferência antes de aparecer. Ao
citar, escreva o nome com inicial maiúscula seguido do número, como em ``a
Figura 3 mostra'', e não ``a figura 3''. O mesmo vale para Tabela, Quadro,
Código, Equação e Seção. Os comandos \verb|\cref| e \verb|\Cref| já escrevem o
nome dessa forma. Não basta citar. Explique no texto o que deve ser observado. Diagramas de
arquitetura, de classes e de sequência pedem explicação no texto, elemento por
elemento ou pelo menos das partes principais.

\textbf{Como obter figuras de boa qualidade}
\begin{itemize}
  \item Prefira figuras \textbf{vetoriais} (PDF, EPS ou desenhadas em TikZ),
        que não perdem nitidez com o zoom nem na impressão. Esta figura foi
        desenhada em TikZ diretamente no texto.
  \item Diagramas UML feitos no draw.io (diagrams.net), no PlantUML, no Astah
        ou em ferramentas semelhantes devem ser exportados em PDF ou SVG
        (convertido para PDF), nunca capturados com \emph{print screen}.
  \item Capturas de tela do sistema devem ser salvas em PNG, na resolução
        original e com o zoom do navegador em 100\%. Recorte apenas a área de
        interesse e evite JPEG, porque a compressão borra textos e linhas.
  \item Ajuste apenas a largura (\verb|width=|). Alterar largura e altura ao
        mesmo tempo distorce a imagem.
  \item O texto dentro da figura deve ter tamanho próximo ao das legendas.
        Se for preciso dar zoom para ler, a figura precisa ser refeita ou
        dividida.
  \item Use cores que continuem distinguíveis em escala de cinza e para
        pessoas com daltonismo.
  \item Mantenha o mesmo estilo e o mesmo alinhamento (centralizado) em todas
        as figuras do trabalho, como pede a NBR 14724:2024.
\end{itemize}
\end{instrucao}

% ------------------------------------------------------------------------------
\section{Usabilidade e a escala SUS}
\label{sec:fundamentacao:usabilidade}

A usabilidade é a medida em que um sistema pode ser usado por usuários
específicos para atingir objetivos com eficácia, eficiência e satisfação em um
contexto de uso \cite{iso2018usability}. Um dos instrumentos mais usados para medir a satisfação é a
\ac{sus}, um questionário de dez afirmações respondidas em escala de cinco
pontos que resulta em uma pontuação de 0 a 100 \cite{brooke1996sus}.
\textcite{bangor2009determining} associaram faixas dessa pontuação a adjetivos,
o que facilita a interpretação. Pontuações em torno de 71 correspondem ao
adjetivo ``bom''. Pontuações em torno de 86 correspondem a ``excelente''.

% ------------------------------------------------------------------------------
\section{Tecnologias utilizadas}
\label{sec:fundamentacao:tecnologias}

O \cref{qua:tecnologias} resume as tecnologias usadas no desenvolvimento, com
a versão e o papel de cada uma. A escolha do Django se deveu à experiência
prévia do autor com Python, ao sistema de autenticação e de administração já
incluídos no arcabouço e à documentação extensa
\cite{django2026documentation}. O PostgreSQL foi escolhido por ser o banco de
dados já mantido pela equipe de tecnologia do centro, o que facilita a
implantação.

\begin{quadro}[htbp]
  \centering
  \caption{Tecnologias utilizadas no desenvolvimento}
  \label{qua:tecnologias}
  \small
  \begin{tabularx}{\textwidth}{|l|c|X|}
    \hline
    \textbf{Tecnologia} & \textbf{Versão} & \textbf{Papel no sistema} \\
    \hline
    Python & 3.12 & Linguagem de programação do servidor \\
    \hline
    Django & 5.2 & Arcabouço web (rotas, visões, \acs{orm}, autenticação) \\
    \hline
    PostgreSQL & 16 & Banco de dados relacional \\
    \hline
    Bootstrap & 5.3 & Componentes de interface e leiaute responsivo \\
    \hline
    pytest & 8.3 & Execução dos testes automatizados \\
    \hline
    Docker & 27 & Empacotamento e implantação em contêineres \\
    \hline
  \end{tabularx}
  \fonteautor
\end{quadro}

% ==============================================================================
% Capítulo 3: Soluções relacionadas
% ==============================================================================
\chapter{Soluções relacionadas}
\label{cap:solucoes}

\begin{instrucao}
\orientacao{O que vai em Soluções Relacionadas}
Em um trabalho tecnológico, este capítulo corresponde aos trabalhos
relacionados de um trabalho científico. As competências do TCC pedem que o
projeto seja relacionado às soluções de mercado equivalentes, demonstrando sua
relevância. Analise:
\begin{itemize}
  \item produtos e serviços que resolvem o mesmo problema ou parte dele;
  \item sistemas institucionais já existentes;
  \item trabalhos acadêmicos (artigos e TCCs) que desenvolveram sistemas
        semelhantes.
\end{itemize}
Defina critérios de comparação ligados aos requisitos do seu sistema e
aplique-os a todas as soluções. Termine com um quadro comparativo e um
parágrafo que mostre o que o seu sistema oferece de diferente. Se alguma
solução já resolve bem o problema, diga isso e explique por que ela não serve
ao seu contexto (custo, integração, privacidade dos dados).

Ao descrever produtos comerciais, informe a data de consulta, porque
funcionalidades e preços mudam. Cite a página do produto como referência
(\verb|@misc| com \verb|url| e \verb|urldate|).
\end{instrucao}

As soluções usadas hoje para organizar monitorias se dividem em ferramentas
genéricas de agendamento, planilhas e formulários compartilhados e o sistema
acadêmico institucional. Esta análise considerou os critérios que emergiram
das entrevistas com os usuários, descritas no \cref{cap:requisitos}.

% ------------------------------------------------------------------------------
\section{Ferramentas genéricas de agendamento}

Serviços de agendamento \emph{on-line} permitem que uma pessoa publique seus
horários livres e que outras reservem esses horários por um endereço web.
Eles evitam reservas simultâneas e enviam lembretes, mas são pensados para
atendimentos individuais de profissionais. Não associam a reserva a uma
disciplina, não registram a presença do estudante nem o assunto atendido e,
nas versões gratuitas, costumam limitar o número de agendas por conta. Além
disso, exigem que os estudantes informem dados pessoais a um serviço externo
à universidade.

% ------------------------------------------------------------------------------
\section{Planilhas e formulários compartilhados}

Planilhas e formulários são a solução mais usada pelos monitores entrevistados.
São gratuitos e conhecidos, mas não impedem que dois estudantes reservem o
mesmo horário, não enviam avisos e dependem do monitor para consolidar os
registros de frequência ao fim do mês. Os professores relataram que os
registros chegam incompletos e em formatos diferentes de um monitor para outro.

% ------------------------------------------------------------------------------
\section{Sistema acadêmico institucional}

O sistema acadêmico da universidade registra os monitores vinculados a cada
disciplina e controla o pagamento de bolsas, mas não oferece funcionalidades
de agendamento ou de registro de atendimentos. Integrar o Monitora a esse
sistema seria desejável, porém exigiria acesso a interfaces de programação
que não estão disponíveis para projetos de estudantes.

% ------------------------------------------------------------------------------
\section{Comparação}

O \cref{qua:comparacao} compara as três categorias com o sistema proposto. O
Monitora é a única opção que reúne agendamento sem conflitos, registro de
frequência e de assuntos atendidos e relatórios por disciplina, mantendo os
dados dos estudantes em infraestrutura da universidade.

\begin{quadro}[htbp]
  \centering
  \caption{Comparação entre as soluções existentes e o Monitora}
  \label{qua:comparacao}
  \small
  \begin{tabularx}{\textwidth}{|X|c|c|c|c|}
    \hline
    \textbf{Critério} & \textbf{Agendadores} & \textbf{Planilhas} &
      \textbf{\begin{tabular}[c]{@{}c@{}}Sistema\\acadêmico\end{tabular}} &
      \textbf{Monitora} \\
    \hline
    Impede reservas simultâneas & Sim & Não & Não & Sim \\
    \hline
    Associa atendimento à disciplina & Não & Parcial & Sim & Sim \\
    \hline
    Registra frequência e assunto & Não & Parcial & Não & Sim \\
    \hline
    Envia lembretes & Sim & Não & Não & Sim \\
    \hline
    Gera relatórios por disciplina & Não & Não & Parcial & Sim \\
    \hline
    Dados em infraestrutura da Ufes & Não & Não & Sim & Sim \\
    \hline
  \end{tabularx}
  \fonte{Elaborado pelo próprio autor com base nas entrevistas e em consultas
    realizadas em outubro de 2026.}
\end{quadro}

\begin{instrucao}
\orientacao{Quadro ou tabela?}
A NBR 14724 distingue os dois. O \textbf{quadro} organiza informações
textuais em linhas e colunas e tem bordas fechadas, com linhas verticais,
como o \cref{qua:comparacao}. A \textbf{tabela} apresenta dados, em geral
numéricos, e segue as normas do IBGE, com bordas laterais abertas e sem linhas
verticais (veja a \cref{tab:cobertura} no \cref{cap:testes}). Quadros e
tabelas têm listas e numeração próprias. A identificação fica acima e a fonte
abaixo, como nas figuras.
\begin{exemplocodigo}
\begin{quadro}[htbp]
  \centering
  \caption{Comparação entre as soluções}
  \label{qua:comparacao}
  \small
  \begin{tabular}{|l|c|c|}
    \hline
    \textbf{Critério} & \textbf{Planilhas} & \textbf{Monitora} \\
    \hline
    Impede reservas simultâneas & Não & Sim \\
    \hline
  \end{tabular}
  \fonteautor
\end{quadro}
\end{exemplocodigo}
Para montar tabelas e quadros mais elaborados, o site
\href{https://www.tablesgenerator.com/}{Tables Generator}
(\url{https://www.tablesgenerator.com/}) gera o código LaTeX a partir de uma
planilha. Depois de colar o código, ajuste-o ao padrão do modelo (ambiente
\verb|table| ou \verb|quadro|, \verb|\caption| acima e \verb|\fonte| abaixo).
\end{instrucao}

% ==============================================================================
% Capítulo 4: Especificação de requisitos
% ==============================================================================
\chapter{Especificação de requisitos}
\label{cap:requisitos}

\begin{instrucao}
\orientacao{O que vai na Especificação de Requisitos}
Este capítulo descreve o que o sistema deve fazer e para quem. Inclua:
\begin{itemize}
  \item como os requisitos foram levantados (entrevistas, questionários,
        observação, análise de documentos) e com quem;
  \item os perfis de usuário e as partes interessadas;
  \item os requisitos funcionais, como casos de uso, histórias de usuário ou
        abordagem equivalente, conforme as competências do TCC;
  \item os requisitos não funcionais, escritos de forma verificável;
  \item as regras de negócio;
  \item os protótipos de interface usados para validar os requisitos.
\end{itemize}
Identifique cada requisito com um código (RF01, RNF01, RN01), porque eles
serão citados no projeto, na implementação e nos testes. Se a especificação
completa for longa, mantenha no texto os requisitos principais e coloque o
documento completo em um apêndice.
\end{instrucao}

Este capítulo apresenta os requisitos do Monitora. A
\cref{sec:requisitos:levantamento} descreve como eles foram levantados, as
\cref{sec:requisitos:funcionais,sec:requisitos:naofuncionais} apresentam os
requisitos funcionais e não funcionais e a \cref{sec:requisitos:prototipo}
mostra o protótipo de interface usado para validá-los.

% ------------------------------------------------------------------------------
\section{Levantamento e perfis de usuário}
\label{sec:requisitos:levantamento}

Os requisitos foram levantados em entrevistas semiestruturadas com quatro
monitores das disciplinas de Programação I e Estrutura de Dados I e com os dois
professores que os supervisionam. As entrevistas duraram cerca de 30 minutos
e seguiram um roteiro com perguntas sobre como os atendimentos são marcados,
registrados e relatados. Foram identificados três perfis de usuário:
\begin{alineas}
  \item \textbf{estudante}: consulta os horários disponíveis, reserva e cancela
        atendimentos;
  \item \textbf{monitor}: publica seus horários, confirma a presença dos
        estudantes e registra o assunto de cada atendimento;
  \item \textbf{professor}: cadastra os monitores da disciplina e consulta os
        relatórios de frequência e de assuntos atendidos.
\end{alineas}

% ------------------------------------------------------------------------------
\section{Requisitos funcionais}
\label{sec:requisitos:funcionais}

Os \acp{rf} foram registrados como histórias de usuário \cite{cohn2004user} e
priorizados com os professores supervisores em alta, média e baixa
prioridade. O \cref{qua:historias} apresenta as histórias de alta e média
prioridade, implementadas nesta versão do sistema.

\begin{quadro}[htbp]
  \centering
  \caption{Histórias de usuário do Monitora}
  \label{qua:historias}
  \small
  \begin{tabularx}{\textwidth}{|c|X|c|}
    \hline
    \textbf{Código} & \textbf{História de usuário} & \textbf{Prioridade} \\
    \hline
    RF01 & Como monitor, quero cadastrar meus horários semanais de atendimento
      para que os estudantes saibam quando estou disponível. & Alta \\
    \hline
    RF02 & Como estudante, quero ver os horários livres de uma disciplina para
      escolher o mais conveniente. & Alta \\
    \hline
    RF03 & Como estudante, quero reservar um horário para garantir meu
      atendimento. & Alta \\
    \hline
    RF04 & Como estudante, quero cancelar uma reserva para liberar o horário
      para outro colega. & Alta \\
    \hline
    RF05 & Como monitor, quero registrar a presença e o assunto de cada
      atendimento para compor meu relatório mensal. & Alta \\
    \hline
    RF06 & Como professor, quero cadastrar os monitores da minha disciplina para
      que eles possam publicar horários. & Média \\
    \hline
    RF07 & Como professor, quero emitir o relatório de frequência de cada
      monitor para comprovar as horas trabalhadas. & Média \\
    \hline
    RF08 & Como estudante, quero receber um lembrete na véspera do atendimento
      para não esquecer o compromisso. & Média \\
    \hline
  \end{tabularx}
  \fonteautor
\end{quadro}

Cada história tem critérios de aceitação que orientaram os testes. Para a
RF03, por exemplo, foram definidos os seguintes critérios:
\begin{alineas}
  \item a reserva só é aceita se o horário estiver livre;
  \item o estudante não pode ter duas reservas no mesmo horário;
  \item após a reserva, o horário deixa de aparecer como livre para os demais
        estudantes.
\end{alineas}

\begin{instrucao}
\orientacao{Histórias de usuário, casos de uso e critérios de aceitação}
Escolha uma forma de especificação e use-a de modo consistente. Histórias de
usuário são adequadas a processos ágeis. Casos de uso, com diagrama e
descrição textual (atores, pré-condições, fluxo principal e fluxos
alternativos), são adequados quando há muitos fluxos alternativos ou regras
complexas. Em ambos os casos, escreva critérios de aceitação verificáveis,
porque eles orientam os testes descritos no \cref{cap:testes}.
\end{instrucao}

% ------------------------------------------------------------------------------
\section{Requisitos não funcionais e regras de negócio}
\label{sec:requisitos:naofuncionais}

Os \acp{rnf} foram definidos de forma que pudessem ser verificados ao final do
desenvolvimento:
\begin{alineas}
  \item \textbf{RNF01 (usabilidade)}: o sistema deve obter pontuação média
        igual ou superior a 70 na escala \acs{sus} em avaliação com usuários;
  \item \textbf{RNF02 (compatibilidade)}: a interface deve funcionar em
        telas a partir de 360 \emph{pixels} de largura, nas versões atuais dos
        navegadores Chrome, Firefox e Safari;
  \item \textbf{RNF03 (desempenho)}: as páginas de consulta de horários devem
        ser carregadas em menos de dois segundos com 100 usuários simultâneos;
  \item \textbf{RNF04 (segurança)}: o acesso deve exigir autenticação e cada
        perfil deve ter acesso apenas às funcionalidades que lhe competem;
  \item \textbf{RNF05 (manutenibilidade)}: o código deve seguir o guia de
        estilo PEP~8 \cite{vanrossum2001pep8} e ter cobertura de testes
        automatizados de pelo menos 80\%.
\end{alineas}

As principais regras de negócio são as seguintes. Um estudante pode ter no
máximo duas reservas ativas por disciplina (RN01). Cancelamentos devem ocorrer
com pelo menos duas horas de antecedência (RN02). Um atendimento não
registrado pelo monitor em até 48 horas é marcado como pendente no relatório
do professor (RN03).

\begin{instrucao}
\orientacao{Requisitos não funcionais verificáveis}
``O sistema deve ser rápido'' ou ``o sistema deve ser fácil de usar'' não são
requisitos verificáveis. Associe cada requisito não funcional a uma medida e
a um valor esperado, como nos exemplos acima. Assim, no \cref{cap:testes}, é
possível mostrar se cada um foi atendido. A norma ISO/IEC 25010 oferece um
vocabulário de atributos de qualidade que ajuda a organizá-los.
\end{instrucao}

% ------------------------------------------------------------------------------
\section{Protótipo de interface}
\label{sec:requisitos:prototipo}

Antes da implementação, foram elaborados protótipos de baixa fidelidade das
telas principais, apresentados aos monitores e professores para validar os
requisitos. A \cref{fig:prototipo} mostra o protótipo da tela de reserva. A
validação levou a duas mudanças. A primeira foi exibir os horários em grade
semanal, e não em lista. A segunda foi mostrar o nome do monitor em cada
horário, porque os estudantes preferem ser atendidos por monitores
específicos.

\begin{figure}[htbp]
  \centering
  \caption{Protótipo de baixa fidelidade da tela de reserva de horários}
  \label{fig:prototipo}
  \begin{tikzpicture}[
      font=\scriptsize,
      livre/.style={draw, fill=white, minimum width=1.7cm, minimum height=7mm, align=center},
      ocupado/.style={draw, fill=gray!30, minimum width=1.7cm, minimum height=7mm, align=center},
    ]
    % janela
    \draw[thick, rounded corners=2pt] (0,0) rectangle (12.4,7.2);
    \fill[gray!15] (0,6.5) rectangle (12.4,7.2);
    \draw (0,6.5) -- (12.4,6.5);
    \node[anchor=west, font=\footnotesize\bfseries] at (0.2,6.85) {Monitora};
    \node[anchor=east] at (12.2,6.85) {Estudante: Ana Souza \quad Sair};
    % menu lateral
    \draw (0,0) rectangle (2.6,6.5);
    \node[anchor=west, font=\scriptsize\bfseries] at (0.2,6.1) {Disciplinas};
    \node[anchor=west, draw, fill=gray!15, minimum width=2.2cm] at (0.2,5.5) {Programação I};
    \node[anchor=west] at (0.2,4.9) {Estrutura de Dados I};
    \node[anchor=west] at (0.2,4.3) {Minhas reservas};
    % grade semanal
    \node[anchor=west, font=\footnotesize\bfseries] at (2.9,6.05) {Horários disponíveis, semana de 10 a 14 de novembro};
    \foreach \d [count=\i] in {Seg, Ter, Qua, Qui, Sex} {
      \node at (1.86+\i*1.88, 5.45) {\textbf{\d}};
    }
    \foreach \h [count=\j] in {14h, 15h, 16h, 17h} {
      \node[anchor=east] at (2.88, 5.4-\j*1.0) {\h};
    }
    \node[livre]   at (3.74,4.4) {Livre\\(João)};
    \node[ocupado] at (5.62,4.4) {Reservado};
    \node[livre]   at (7.50,4.4) {Livre\\(Maria)};
    \node[livre]   at (3.74,3.4) {Livre\\(João)};
    \node[livre]   at (9.38,3.4) {Livre\\(Pedro)};
    \node[ocupado] at (7.50,2.4) {Reservado};
    \node[livre]   at (11.26,2.4) {Livre\\(Maria)};
    \node[livre]   at (5.62,1.4) {Livre\\(Pedro)};
    % legenda e botão
    \node[draw, fill=white, minimum width=4mm, minimum height=3mm] at (3.3,0.45) {};
    \node[anchor=west] at (3.55,0.45) {livre};
    \node[draw, fill=gray!30, minimum width=4mm, minimum height=3mm] at (4.6,0.45) {};
    \node[anchor=west] at (4.85,0.45) {reservado};
    \node[draw, thick, rounded corners=2pt, fill=gray!15, minimum width=2.4cm, minimum height=6mm] at (10.9,0.45) {Reservar horário};
  \end{tikzpicture}
  \fonteautor
\end{figure}

\begin{instrucao}
\orientacao{Protótipos}
Protótipos ajudam a validar requisitos antes de implementar e mitigam o risco
de construir algo que os usuários não querem, como pedem as competências do
TCC. Podem ser esboços em papel fotografados, telas feitas no Figma ou no
Penpot (exportadas em PDF ou SVG) ou desenhos simples em TikZ, como a
\cref{fig:prototipo}. Relate o que mudou nos requisitos a partir da validação
com os usuários.
\end{instrucao}

% ==============================================================================
% Capítulo 5: Projeto do sistema
% ==============================================================================
\chapter{Projeto do sistema}
\label{cap:projeto}

\begin{instrucao}
\orientacao{O que vai no Projeto do Sistema}
Este capítulo mostra como o sistema foi organizado para atender aos
requisitos. As competências do TCC pedem a definição de uma arquitetura
adequada, modelos de estrutura e de comportamento e a análise do projeto sob a
perspectiva de atributos de qualidade. Inclua:
\begin{itemize}
  \item a \textbf{arquitetura}: componentes, camadas, serviços externos e como
        se comunicam;
  \item o \textbf{modelo estrutural}: diagrama de classes ou modelo de dados;
  \item o \textbf{modelo comportamental}: diagramas de sequência, de atividades
        ou de estados para os fluxos mais importantes;
  \item as \textbf{decisões de projeto}: alternativas consideradas e por que a
        escolhida foi adotada, relacionando-as aos requisitos.
\end{itemize}
Não inclua diagramas sem explicação. Cada diagrama deve ser apresentado no
texto, que destaca o que ele mostra de relevante. Diagramas grandes demais
para uma página devem ser divididos ou ir para um apêndice.
\end{instrucao}

Este capítulo apresenta o projeto do Monitora. A
\cref{sec:projeto:arquitetura} descreve a arquitetura, a
\cref{sec:projeto:modelo} apresenta o modelo de classes e a
\cref{sec:projeto:decisoes} discute as principais decisões de projeto.

% ------------------------------------------------------------------------------
\section{Arquitetura}
\label{sec:projeto:arquitetura}

O sistema segue uma arquitetura em camadas, apresentada na
\cref{fig:arquitetura}. O servidor web Nginx recebe as requisições e as
encaminha à aplicação Django, executada pelo servidor de aplicação Gunicorn. A
aplicação está dividida em três módulos, que correspondem aos perfis de
usuário e aos grupos de requisitos. O módulo \texttt{contas} trata da
autenticação e dos perfis, o módulo \texttt{agenda} trata dos horários e das
reservas (RF01 a RF05) e o módulo \texttt{relatorios} gera os relatórios dos
professores (RF07). Os lembretes (RF08) são enviados por uma tarefa agendada,
executada fora do ciclo das requisições para não atrasar as respostas.

\begin{figure}[htbp]
  \centering
  \caption{Arquitetura do Monitora}
  \label{fig:arquitetura}
  \begin{tikzpicture}[
      font=\footnotesize,
      comp/.style={draw, rounded corners=2pt, minimum height=8mm,
                   minimum width=2.6cm, align=center, fill=white},
      seta/.style={-{Stealth}, thick},
    ]
    \node[comp] (nav) {Navegador};
    \node[comp, below=7mm of nav] (nginx) {Nginx};
    \node[comp, below=7mm of nginx] (gunicorn) {Gunicorn};
    \node[comp, below=10mm of gunicorn, xshift=-3.2cm] (contas) {\texttt{contas}};
    \node[comp, below=10mm of gunicorn] (agenda) {\texttt{agenda}};
    \node[comp, below=10mm of gunicorn, xshift=3.2cm] (relat) {\texttt{relatorios}};
    \node[comp, below=12mm of agenda, fill=gray!10] (bd) {PostgreSQL};
    \node[comp, right=12mm of relat] (tarefa) {Tarefa agendada\\(lembretes)};
    \begin{scope}[on background layer]
      \node[draw, dashed, rounded corners=3pt, fill=gray!5, inner sep=3mm,
            fit=(contas)(agenda)(relat),
            label={[anchor=north west, font=\scriptsize]north west:Aplicação Django}] (django) {};
      \node[draw, dotted, rounded corners=3pt, inner sep=4mm,
            fit=(nginx)(django)(bd)(tarefa),
            label={[anchor=south east, font=\scriptsize]south east:Contêineres Docker}] {};
    \end{scope}
    \draw[seta] (nav) -- node[right, font=\scriptsize] {HTTPS} (nginx);
    \draw[seta] (nginx) -- (gunicorn);
    \draw[seta] (gunicorn) -- (django.north);
    \draw[seta] (contas) |- (bd);
    \draw[seta] (agenda) -- (bd);
    \draw[seta] (relat) |- (bd);
    \draw[seta] (tarefa.south) |- ([yshift=-2.5mm]bd.east);
    \draw[seta] (tarefa.east) -- ++(9mm,0) node[right, align=left, font=\scriptsize] {servidor\\de e-mail};
  \end{tikzpicture}
  \fonteautor
\end{figure}

% ------------------------------------------------------------------------------
\section{Modelo de classes}
\label{sec:projeto:modelo}

A \cref{fig:classes} apresenta o diagrama de classes de domínio, em notação
\ac{uml}, com os atributos principais. Uma \texttt{Disciplina} tem vários
monitores e cada \texttt{Monitor} publica \texttt{Horarios} semanais. Cada
ocorrência de um horário em uma data específica pode receber no máximo uma
\texttt{Reserva}, feita por um \texttt{Estudante}. Quando o atendimento
acontece, o monitor registra o assunto e a presença na própria reserva, o que
evita uma entidade adicional e simplifica os relatórios. As classes
\texttt{Monitor}, \texttt{Estudante} e \texttt{Professor} são perfis de um
mesmo usuário do sistema de autenticação do Django, o que permite que um
monitor também reserve atendimentos em outra disciplina.

\begin{figure}[htbp]
  \centering
  \caption{Diagrama de classes de domínio do Monitora}
  \label{fig:classes}
  \begin{tikzpicture}[
      font=\scriptsize,
      classe/.style={draw, rectangle split, rectangle split parts=2,
                     rectangle split part fill={gray!15,white},
                     text width=3.3cm, align=left, inner sep=3pt},
      assoc/.style={thick},
    ]
    \node[classe] (disc) {\centering\textbf{Disciplina}
      \nodepart{second} codigo: str\\ nome: str\\ semestre: str};
    \node[classe, right=2.2cm of disc] (mon) {\centering\textbf{Monitor}
      \nodepart{second} usuario: Usuario\\ disciplina: Disciplina};
    \node[classe, below=1.6cm of mon] (hor) {\centering\textbf{Horario}
      \nodepart{second} dia\_semana: int\\ inicio: time\\ fim: time};
    \node[classe, below=1.6cm of hor] (res) {\centering\textbf{Reserva}
      \nodepart{second} data: date\\ assunto: str\\ presente: bool};
    \node[classe, left=2.2cm of res] (est) {\centering\textbf{Estudante}
      \nodepart{second} usuario: Usuario\\ matricula: str};
    \node[classe, below=1.6cm of disc] (prof) {\centering\textbf{Professor}
      \nodepart{second} usuario: Usuario\\ siape: str};
    \draw[assoc] (disc) -- node[above, pos=0.15] {1} node[above, pos=0.85] {1..*} (mon);
    \draw[assoc] (prof) -- node[left, pos=0.2] {1..*} node[left, pos=0.8] {1..*} (disc);
    \draw[assoc] (mon) -- node[right, pos=0.2] {1} node[right, pos=0.8] {0..*} (hor);
    \draw[assoc] (hor) -- node[right, pos=0.2] {1} node[right, pos=0.8] {0..*} (res);
    \draw[assoc] (est) -- node[above, pos=0.15] {1} node[above, pos=0.85] {0..*} (res);
  \end{tikzpicture}
  \fonteautor
\end{figure}

% ------------------------------------------------------------------------------
\section{Decisões de projeto e atributos de qualidade}
\label{sec:projeto:decisoes}

A decisão mais importante do projeto foi como garantir que um horário não seja
reservado por dois estudantes ao mesmo tempo (RF03). A primeira alternativa
considerada foi verificar a disponibilidade na visão antes de gravar a
reserva. Essa verificação, porém, não impede que duas requisições simultâneas
passem pela checagem antes de qualquer uma gravar. Por isso, adotou-se uma
restrição de unicidade no banco de dados sobre o par (horário, data), de modo
que a segunda gravação falhe e a aplicação informe ao estudante que o horário
acabou de ser ocupado. O \cref{cod:conflito}, no \cref{cap:implementacao},
mostra a implementação dessa regra.

Do ponto de vista dos atributos de qualidade, a separação em módulos favorece
a manutenibilidade (RNF05), porque cada módulo pode ser testado de forma
isolada. A execução dos lembretes em tarefa separada contribui para o
desempenho (RNF03). O uso do sistema de autenticação e de permissões do
próprio Django reduz o risco de falhas de segurança em código escrito do zero
(RNF04).

\begin{instrucao}
\orientacao{Decisões de projeto}
Bancas valorizam o raciocínio por trás das escolhas. Para cada decisão
relevante, apresente o problema, as alternativas consideradas, a escolhida e o
motivo, relacionando-a aos requisitos e aos atributos de qualidade. Diagramas
de sequência ajudam a explicar fluxos com várias etapas, como o de reserva
descrito acima. Se o diagrama ficar grande, mostre apenas o fluxo principal.
\end{instrucao}

% ==============================================================================
% Capítulo 6: Implementação
% ==============================================================================
\chapter{Implementação}
\label{cap:implementacao}

\begin{instrucao}
\orientacao{O que vai na Implementação}
Este capítulo mostra como o projeto foi transformado em software. Inclua:
\begin{itemize}
  \item o ambiente e as ferramentas de desenvolvimento;
  \item a organização do código e os padrões de documentação e codificação
        adotados;
  \item trechos de código que mostrem soluções relevantes ou não óbvias, e
        não o código inteiro;
  \item as telas principais do sistema, relacionadas aos requisitos que
        atendem;
  \item como o sistema foi implantado.
\end{itemize}
O texto não deve ser um manual do usuário. Mostre as telas para evidenciar que
os requisitos foram atendidos e explique o que cada uma permite fazer. Um
manual completo, se existir, vai para um apêndice. Informe o endereço do
repositório com o código-fonte.
\end{instrucao}

Este capítulo descreve a implementação do Monitora. A
\cref{sec:implementacao:organizacao} apresenta a organização do código, a
\cref{sec:implementacao:reserva} detalha a implementação da reserva de
horários, a \cref{sec:implementacao:implantacao} descreve a implantação e a
\cref{sec:implementacao:ia} relata o uso de ferramentas de inteligência
artificial.

% ------------------------------------------------------------------------------
\section{Ambiente e organização do código}
\label{sec:implementacao:organizacao}

O sistema foi desenvolvido no Visual Studio Code, com controle de versões no
Git e repositório hospedado no GitHub\footnote{O endereço do repositório deve
ser informado aqui, por exemplo, \url{https://github.com/usuario/monitora}.}.
O código segue a estrutura padrão de projetos Django, com um diretório por
módulo (\texttt{contas}, \texttt{agenda} e \texttt{relatorios}), cada um com
seus modelos, visões, \emph{templates} e testes. A formatação segue o guia
PEP~8 \cite{vanrossum2001pep8}, verificada automaticamente pela ferramenta
Ruff a cada \emph{commit}. Funções e classes públicas são documentadas com
\emph{docstrings}.

% ------------------------------------------------------------------------------
\section{Reserva de horários}
\label{sec:implementacao:reserva}

O \cref{cod:conflito} mostra o modelo \texttt{Reserva} e a função que cria
uma reserva. A restrição de unicidade declarada nas linhas 13 a 17 implementa
a decisão de projeto discutida na \cref{sec:projeto:decisoes}. Se duas
requisições tentarem reservar o mesmo horário na mesma data, o banco de dados
aceita apenas a primeira e a segunda gera um \texttt{IntegrityError}. A
função \texttt{reservar} (linhas 24 a 33) converte esse erro em uma exceção
do domínio, tratada pela visão, que exibe ao estudante a mensagem de que o
horário acabou de ser ocupado.

\begin{codigo}[htbp]
  \caption{Modelo \texttt{Reserva} e função de reserva com tratamento de conflito}
  \label{cod:conflito}
  \lstinputlisting[language=Python]{codigos/reservas.py}
  \fonteautor
\end{codigo}

\begin{instrucao}
\orientacao{Códigos-fonte}
Inclua apenas trechos curtos e relevantes, em geral com menos de 30 linhas.
O código completo fica no repositório. Coloque o trecho no ambiente
\verb|codigo|, que tem legenda acima, fonte abaixo e lista própria (Lista de
códigos). Para não copiar o código para dentro do \verb|.tex|, guarde o
arquivo na pasta \verb|codigos/| e use \verb|\lstinputlisting|. As opções
\verb|firstline| e \verb|lastline| permitem mostrar só parte do arquivo.
\begin{exemplocodigo}
\begin{codigo}[htbp]
  \caption{Função de reserva}
  \label{cod:reserva}
  \lstinputlisting[language=Python, firstline=24, lastline=33]{codigos/reservas.py}
  \fonteautor
\end{codigo}
\end{exemplocodigo}
Para trechos muito curtos, escritos direto no texto, use o ambiente
\verb|lstlisting| dentro do \verb|codigo|. Explique o código no texto citando
os números das linhas, como no parágrafo acima. Nunca use capturas de tela de
código, que não podem ser lidas por leitores de tela nem copiadas.

Algoritmos devem ser apresentados em pseudocódigo, e não em uma linguagem
específica. O modelo científico mostra como fazer isso com o ambiente
\verb|algoritmo| e o pacote \verb|algpseudocode|.
\end{instrucao}

% ------------------------------------------------------------------------------
\section{Telas do sistema}
\label{sec:implementacao:telas}

\begin{instrucao}
\orientacao{Telas}
Mostre as telas que correspondem aos requisitos principais, uma por figura,
com legenda que diga o que a tela permite fazer, e cite no texto os
requisitos atendidos. Capture as telas em PNG, na resolução original e com
zoom de 100\%, preenchidas com dados fictícios (nunca dados pessoais reais de
estudantes). Para incluir uma captura de tela:
\begin{exemplocodigo}
\begin{figure}[htbp]
  \centering
  \caption{Tela de reserva de horários (RF02 e RF03)}
  \label{fig:tela-reserva}
  \includegraphics[width=0.9\textwidth]{figuras/tecnologico/tela-reserva.png}
  \fonteautor
\end{figure}
\end{exemplocodigo}
Duas telas relacionadas podem ser apresentadas lado a lado com o pacote
\verb|subcaption|, usando \verb|\begin{subfigure}{0.48\textwidth}| dentro do
ambiente \verb|figure|.
\end{instrucao}

A tela de reserva segue o protótipo apresentado na \cref{fig:prototipo}. O
estudante escolhe a disciplina no menu lateral e vê a grade semanal com os
horários livres e o nome de cada monitor (RF02). Ao selecionar um horário, uma
janela de confirmação mostra a data, o monitor e o local do atendimento
(RF03). As reservas do estudante ficam disponíveis no item ``Minhas
reservas'', no qual também é possível cancelá-las (RF04).

% ------------------------------------------------------------------------------
\section{Implantação}
\label{sec:implementacao:implantacao}

O sistema é distribuído em três contêineres Docker, descritos em um arquivo
do Docker Compose, que executam o Nginx, a aplicação Django com o Gunicorn e
o PostgreSQL. A tarefa de envio de lembretes é executada
diariamente pelo agendador do próprio servidor. A versão avaliada no
\cref{cap:testes} foi implantada em uma máquina virtual de homologação
cedida pelo departamento.

% ------------------------------------------------------------------------------
\section{Uso de ferramentas de inteligência artificial}
\label{sec:implementacao:ia}

Durante a implementação, o GitHub Copilot \cite{github2026copilot}, usado
como extensão do Visual Studio Code com o modelo padrão do plano gratuito para
estudantes, sugeriu complementos de código em funções de formatação dos
relatórios e em testes automatizados. Todas as sugestões aceitas foram
revisadas, ajustadas pelo autor e cobertas por testes. Nenhum dado pessoal,
nenhuma credencial e nenhum dado dos usuários da avaliação foram expostos à
ferramenta. A declaração completa e o detalhamento das interações estão no
\cref{ap:ia}.

\begin{instrucao}
\orientacao{Declaração de uso de inteligência artificial}
O Art. 19 da Portaria Normativa nº 310/2026 da Ufes \cite{ufes2026portaria}
torna obrigatória a declaração do uso de \ac{ia} em qualquer produção
acadêmica, exceto os usos meramente acessórios definidos pela comissão de
governança de IA da Ufes (CPGIA). Assistentes de programação, como o GitHub
Copilot, também precisam ser declarados. Para trabalhos acadêmicos, a portaria
pede três coisas. O modelo atende a cada uma delas em um lugar:
\begin{enumerate}
  \item uma declaração com ferramenta, versão, finalidade e forma de
        contribuição, em seção ou apêndice específico (§ 1º): o
        \cref{ap:ia}, que fica como primeiro apêndice para ser encontrado com
        facilidade;
  \item o detalhamento das interações, das motivações e dos métodos de uso
        (§ 2º): também no \cref{ap:ia};
  \item a citação da ferramenta conforme as normas de citação e a indicação,
        na descrição dos métodos, da versão, dos parâmetros e das bases de
        dados usadas (§ 3º): esta seção, com a ferramenta nas Referências.
\end{enumerate}
A portaria também proíbe atribuir autoria ou coautoria a uma ferramenta de IA
(Art. 17) e inserir nela dados pessoais, credenciais institucionais ou dados
confidenciais. Você continua responsável pela correção e pela originalidade
de tudo o que entregar, inclusive do código (Art. 18). Se não usou nenhuma
ferramenta de IA, apague esta seção e o \cref{ap:ia}. Quando o uso decorrer de
tecnologia assistiva, a declaração pode ser apresentada ao setor de apoio à
inclusão, e não no corpo do trabalho (§ 4º). O modelo oficial de declaração
ainda será publicado pela CPGIA. Quando isso acontecer, ajuste o apêndice a
ele.
\end{instrucao}

% ==============================================================================
% Capítulo 7: Verificação, validação e testes
% ==============================================================================
\chapter{Verificação, validação e testes}
\label{cap:testes}

\begin{instrucao}
\orientacao{O que vai em Verificação, Validação e Testes}
Este capítulo mostra que o sistema funciona e que atende aos requisitos.
Distinga dois aspectos:
\begin{itemize}
  \item \textbf{verificação}: o sistema foi construído corretamente? Testes
        automatizados (unidade, integração, sistema), cobertura, análise
        estática;
  \item \textbf{validação}: o sistema certo foi construído? Avaliação com
        usuários reais, testes de aceitação, medidas de usabilidade.
\end{itemize}
Retome os requisitos funcionais e não funcionais e mostre, para cada um, como
foi verificado e com que resultado. Requisitos não atendidos também devem ser
relatados, com a explicação do motivo.

\textbf{Atenção:} todos os números deste capítulo de exemplo são fictícios.
\end{instrucao}

Este capítulo apresenta a verificação e a validação do Monitora. A
\cref{sec:testes:automatizados} descreve os testes automatizados, a
\cref{sec:testes:rnf} mostra a verificação dos requisitos não funcionais e a
\cref{sec:testes:usuarios} relata a avaliação com usuários.

% ------------------------------------------------------------------------------
\section{Testes automatizados}
\label{sec:testes:automatizados}

Foram escritos 142 testes automatizados com o pytest, divididos em testes de
unidade dos modelos e das regras de negócio, testes de integração das visões
e testes de sistema que simulam o uso pelo navegador. Os critérios de
aceitação das histórias de usuário (\cref{sec:requisitos:funcionais}) foram
convertidos em testes de sistema, de modo que cada história implementada tem
pelo menos um teste correspondente. Os testes são executados automaticamente
a cada envio de código ao repositório.

A \cref{tab:cobertura} apresenta a cobertura de código por módulo. A cobertura
total de 87\% atende ao RNF05. O módulo \texttt{relatorios} tem a menor
cobertura porque parte de seu código formata arquivos para impressão, cujo
resultado foi conferido manualmente.

\begin{table}[htbp]
  \centering
  \caption{Cobertura dos testes automatizados por módulo}
  \label{tab:cobertura}
  \small
  \begin{tabular}{@{}lrrr@{}}
    \toprule
    Módulo & Instruções & Instruções cobertas & Cobertura (\%) \\
    \midrule
    \texttt{contas}     &   320 &   291 & 90,9 \\
    \texttt{agenda}     &   610 &   543 & 89,0 \\
    \texttt{relatorios} &   270 &   211 & 78,1 \\
    \midrule
    Total               & 1.200 & 1.045 & 87,1 \\
    \bottomrule
  \end{tabular}
  \fonte{Elaborado pelo próprio autor a partir do relatório da ferramenta
    coverage.py.}
\end{table}

\begin{instrucao}
\orientacao{Tabelas}
Tabelas seguem as normas de apresentação tabular do IBGE, adotadas pela NBR
14724:
\begin{itemize}
  \item sem linhas verticais e com as laterais abertas;
  \item traços horizontais apenas no topo, abaixo do cabeçalho e no fim (use
        \verb|\toprule|, \verb|\midrule| e \verb|\bottomrule| do pacote
        \verb|booktabs|);
  \item números alinhados à direita, com vírgula decimal, ponto como separador
        de milhar (1.200) e a mesma quantidade de casas decimais em toda a
        coluna;
  \item unidade de medida no título ou no cabeçalho da coluna;
  \item fonte obrigatória abaixo da tabela e notas, se houver, logo depois
        (comandos \verb|\fonte| e \verb|\nota|).
\end{itemize}
\begin{exemplocodigo}
\begin{table}[htbp]
  \centering
  \caption{Cobertura dos testes por módulo}
  \label{tab:cobertura}
  \small
  \begin{tabular}{@{}lrr@{}}
    \toprule
    Módulo & Instruções & Cobertura (\%) \\
    \midrule
    contas &   320 & 90,9 \\
    agenda &   610 & 89,0 \\
    \bottomrule
  \end{tabular}
  \fonteautor
\end{table}
\end{exemplocodigo}
O \href{https://www.tablesgenerator.com/}{Tables Generator}
(\url{https://www.tablesgenerator.com/}) ajuda a montar o código a partir de
uma planilha. Marque a opção \emph{booktabs} e, depois de colar o código,
acrescente \verb|\label| e \verb|\fonte|.
\end{instrucao}

% ------------------------------------------------------------------------------
\section{Verificação dos requisitos não funcionais}
\label{sec:testes:rnf}

O \cref{qua:rnf} resume como cada requisito não funcional foi verificado e o
resultado obtido. Todos foram atendidos. O teste de carga do RNF03 foi
executado com a ferramenta Locust, simulando 100 usuários que consultavam
horários e faziam reservas durante dez minutos.

\begin{quadro}[htbp]
  \centering
  \caption{Verificação dos requisitos não funcionais}
  \label{qua:rnf}
  \small
  \begin{tabularx}{\textwidth}{|c|X|X|c|}
    \hline
    \textbf{Requisito} & \textbf{Método de verificação} & \textbf{Resultado} &
      \textbf{Atendido} \\
    \hline
    RNF01 & Questionário \acs{sus} com 18 usuários & Média de 84,6 & Sim \\
    \hline
    RNF02 & Testes manuais em três navegadores e em telas de 360 a 1.920
      \emph{pixels} & Sem falhas de leiaute & Sim \\
    \hline
    RNF03 & Teste de carga com 100 usuários simultâneos & Tempo médio de
      resposta de 0,8 s & Sim \\
    \hline
    RNF04 & Testes automatizados de permissão por perfil & 24 testes
      aprovados & Sim \\
    \hline
    RNF05 & Ruff e relatório de cobertura & Nenhum alerta e 87\% de cobertura
      & Sim \\
    \hline
  \end{tabularx}
  \fonteautor
\end{quadro}

% ------------------------------------------------------------------------------
\section{Avaliação com usuários}
\label{sec:testes:usuarios}

A avaliação contou com 18 participantes voluntários, sendo 12 estudantes, 4
monitores e 2 professores. Cada participante realizou um roteiro de tarefas
compatível com seu perfil, como reservar e cancelar um horário (estudantes),
registrar a presença em um atendimento (monitores) e emitir um relatório
(professores). Ao final, respondeu ao questionário \acs{sus}
\cite{brooke1996sus}. A pontuação de cada participante foi calculada pela
\cref{eq:sus}, em que $x_i$ é a resposta à afirmação $i$, de 1 (discordo
totalmente) a 5 (concordo totalmente).
\begin{equation}
  \mathit{SUS} = 2,5 \times \left[ \sum_{i \in \{1, 3, 5, 7, 9\}} (x_i - 1)
    + \sum_{i \in \{2, 4, 6, 8, 10\}} (5 - x_i) \right]
  \label{eq:sus}
\end{equation}

A \cref{fig:sus} mostra a pontuação média por perfil. A média geral foi de
84,6, acima da meta de 70 definida no RNF01 e próxima da faixa associada ao
adjetivo ``excelente'' por \textcite{bangor2009determining}. Os professores
deram as menores notas, o que se explica pelos comentários registrados ao fim
das sessões. Os dois sentiram falta de exportar os relatórios em planilha,
funcionalidade que ficou para versões futuras.

\begin{figure}[htbp]
  \centering
  \caption{Pontuação média na escala \acs{sus} por perfil de usuário}
  \label{fig:sus}
  \begin{tikzpicture}
    \begin{axis}[
        ybar, bar width=12mm,
        width=0.7\textwidth, height=5.5cm,
        ymin=0, ymax=100, ylabel={Pontuação \acs{sus}},
        symbolic x coords={Estudantes, Monitores, Professores},
        xtick=data, enlarge x limits=0.25,
        nodes near coords,
        every node near coord/.append style={font=\footnotesize,
          /pgf/number format/fixed, /pgf/number format/fixed zerofill,
          /pgf/number format/precision=1},
        tick label style={font=\footnotesize}, label style={font=\footnotesize},
        axis x line*=bottom, axis y line*=left, clip=false,
      ]
      \addplot[fill=gray!40, draw=black] coordinates {
        (Estudantes,85.4) (Monitores,84.4) (Professores,80.0)
      };
      \draw[dashed] (axis cs:Estudantes,70) -- (axis cs:Professores,70)
        node[right, font=\scriptsize] {meta (70)};
    \end{axis}
  \end{tikzpicture}
  \fonte{Elaborado pelo próprio autor. Participaram 12 estudantes, 4
    monitores e 2 professores.}
\end{figure}

\begin{instrucao}
\orientacao{Pesquisas com pessoas}
Avaliações com usuários envolvem seres humanos e exigem cuidados éticos.
Converse com o orientador, antes de aplicar questionários ou entrevistas,
sobre a necessidade de submeter o projeto a um Comitê de Ética em Pesquisa.
Pesquisas de opinião com participantes não identificados podem ser
dispensadas desse registro, conforme a Resolução CNS nº 510/2016, mas a
decisão deve ser tomada caso a caso. Em qualquer situação, explique aos
participantes o objetivo da avaliação, obtenha o consentimento deles, não
colete dados pessoais desnecessários e apresente os resultados de forma
agregada. Coloque o roteiro de tarefas e o questionário em apêndice.
\end{instrucao}

% ==============================================================================
% Capítulo 8: Considerações finais
% ==============================================================================
\chapter{Considerações finais}
\label{cap:conclusao}

\begin{instrucao}
\orientacao{O que vai nas Considerações Finais}
Retome o problema e os objetivos, mostre como cada objetivo específico foi
atingido e em que capítulo, resuma o que foi entregue, reconheça as
limitações, relate as lições aprendidas e as dificuldades enfrentadas e aponte
trabalhos futuros. Como referência, de duas a quatro páginas costumam ser
suficientes.

Não apresente resultados novos. Em trabalhos tecnológicos, vale informar a
situação do sistema ao final do trabalho (em uso, em homologação, aguardando
implantação) e o que falta para colocá-lo em produção.
\end{instrucao}

Este trabalho apresentou o Monitora, um sistema web para agendamento de
monitorias que substitui o processo informal baseado em mensagens e planilhas.
O sistema permite que monitores publiquem horários, que estudantes reservem
atendimentos e que professores acompanhem a frequência e os assuntos
atendidos, mantendo os dados em infraestrutura da universidade.

Os objetivos específicos foram atingidos da seguinte forma. Os requisitos
foram levantados com monitores e professores e especificados no
\cref{cap:requisitos}. A comparação com as soluções existentes, apresentada no
\cref{cap:solucoes}, mostrou que nenhuma delas atendia ao conjunto de
necessidades identificado. A arquitetura e o modelo de dados foram descritos
no \cref{cap:projeto}. A implementação e a implantação em ambiente de
homologação foram apresentadas no \cref{cap:implementacao}. Por fim, o \cref{cap:testes} mostrou
que todos os requisitos não funcionais foram atendidos e que os usuários
avaliaram positivamente a usabilidade do sistema.

% ------------------------------------------------------------------------------
\section{Limitações e lições aprendidas}
\label{sec:conclusao:limitacoes}

O sistema ainda não está integrado ao sistema acadêmico da universidade, o que
obriga o professor a cadastrar manualmente os monitores a cada semestre. A
avaliação com usuários envolveu um número pequeno de participantes de apenas
duas disciplinas. Por isso, os resultados podem não se repetir em turmas maiores.

Do ponto de vista do processo, a validação dos protótipos antes da
implementação evitou retrabalho, já que duas mudanças importantes na tela de
reserva surgiram nessa etapa. A principal dificuldade foi estimar o esforço das
histórias nos primeiros ciclos, que ficaram sobrecarregados. A partir do
terceiro ciclo, o uso da velocidade observada nos ciclos anteriores tornou o
planejamento mais preciso.

% ------------------------------------------------------------------------------
\section{Trabalhos futuros}
\label{sec:conclusao:futuros}

\begin{instrucao}
\orientacao{Trabalhos futuros}
Evite listas genéricas como ``melhorar a interface''. Para cada sugestão,
explique por que ela é relevante e como poderia ser feita, com base no que o
trabalho mostrou (requisitos de baixa prioridade, comentários dos usuários,
limitações encontradas).
\end{instrucao}

A sugestão mais pedida pelos professores foi exportar os relatórios em
planilha, o que pode ser feito com pouco esforço a partir das consultas já
implementadas. A integração com o sistema acadêmico eliminaria o cadastro
manual de monitores, mas depende da disponibilização de uma interface de
programação pela equipe de tecnologia da universidade. Por fim, um painel com
os assuntos mais procurados em cada disciplina poderia ajudar os professores a
identificar conteúdos que precisam de reforço em sala de aula.


% ------------------------------------------------------------------------------
% Elementos pós-textuais
% ------------------------------------------------------------------------------
\postextual

% Referências (NBR 6023), geradas pelo Biber a partir de referencias.bib.
\printbibliography[heading=bibintoc]

% Apêndices: material elaborado pelo próprio autor.
\begin{apendicesenv}
% ==============================================================================
% Apêndice A: Declaração de uso de inteligência artificial
% Exigida pelo Art. 19 da Portaria Normativa nº 310/2026 da Ufes.
% Se você não usou ferramentas de IA, apague este arquivo, a linha que o
% inclui no arquivo principal e a Seção 6.5 da Implementação.
% ==============================================================================
\chapter{Declaração de uso de inteligência artificial}
\label{ap:ia}

\begin{instrucao}
\orientacao{Por que um apêndice?}
O Art. 19 da Portaria Normativa nº 310/2026 da Ufes \cite{ufes2026portaria}
admite que a declaração fique em seção ou apêndice específico. O apêndice foi
escolhido porque a NBR 14724 não prevê um elemento pré-textual para esse fim e
porque o detalhamento das interações exigido pelo § 2º pode ser longo. Por ser
o primeiro apêndice, ele aparece logo após as Referências e é fácil de
localizar. Na Implementação, uma seção curta cita a ferramenta e remete a este
apêndice.

\orientacao{Como preencher}
\begin{itemize}
  \item No \cref{qua:ia:declaracao}, liste cada ferramenta usada, com a versão
        ou o modelo, a finalidade e a forma de contribuição, indicando as
        partes do trabalho afetadas.
  \item Na seção seguinte, detalhe as interações, informando o que foi pedido,
        por que, como o resultado foi verificado e o que foi aproveitado ou
        descartado.
        Não é preciso transcrever todas as conversas. Descreva os tipos de
        interação com exemplos representativos.
  \item Cite cada ferramenta e inclua-a no arquivo \verb|referencias.bib|
        (veja as entradas \verb|openai2026chatgpt| e \verb|github2026copilot|).
  \item A declaração é escrita na primeira pessoa do singular (``Declaro
        que\ldots''), por ser um compromisso pessoal do autor. O detalhamento
        das interações volta à terceira pessoa, como o restante do trabalho.
  \item Seja específico e honesto. Usos ``meramente acessórios'', que a
        portaria dispensa de declaração, serão definidos pela CPGIA. Na dúvida,
        declare.
\end{itemize}
Substitua o texto de exemplo abaixo pelas ferramentas, versões e usos reais
do seu trabalho. Quando a CPGIA publicar o modelo oficial de declaração,
adapte este apêndice a ele.
\end{instrucao}

% ------------------------------------------------------------------------------
\section{Declaração}

Declaro, em cumprimento ao Art. 19 da Portaria Normativa nº 310/2026 da
\imprimiruniversidade{} \cite{ufes2026portaria}, que utilizei ferramentas de
\acl{ia} na elaboração deste trabalho, nos termos descritos no
\cref{qua:ia:declaracao}. Declaro também que revisei todo o conteúdo gerado
com o auxílio dessas ferramentas, que assumo integral responsabilidade pela
precisão e pela originalidade do texto, do código e dos resultados
apresentados e que nenhuma ferramenta foi tratada como autora ou coautora do
trabalho.

\begin{quadro}[htbp]
  \centering
  \caption{Ferramentas de \acs{ia} utilizadas no trabalho}
  \label{qua:ia:declaracao}
  \small
  \begin{tabularx}{\textwidth}{|>{\raggedright\arraybackslash}p{2.8cm}|>{\raggedright\arraybackslash}X|>{\raggedright\arraybackslash}X|>{\raggedright\arraybackslash}p{2.6cm}|}
    \hline
    \textbf{Ferramenta e versão} & \textbf{Finalidade} &
      \textbf{Forma de contribuição} & \textbf{Partes afetadas} \\
    \hline
    GitHub Copilot, extensão do Visual Studio Code, modelo padrão do plano
      para estudantes \cite{github2026copilot} &
      Apoio à programação &
      Sugestões de complemento de código em funções de formatação de
      relatórios e em testes, revisadas e ajustadas pelo autor &
      Módulo \texttt{relatorios} e testes automatizados
      (\cref{sec:implementacao:organizacao,sec:testes:automatizados}) \\
    \hline
  \end{tabularx}
  \fonteautor
\end{quadro}

% ------------------------------------------------------------------------------
\section{Detalhamento das interações}

\textbf{Complemento de código.} A extensão ficou ativa no editor durante a
implementação do módulo \texttt{relatorios} e dos testes automatizados. Ela
sugeria complementos de linhas e de funções a partir do código já escrito e
dos comentários. A motivação foi reduzir o tempo gasto com código repetitivo,
como a formatação de datas e a montagem de tabelas dos relatórios. As
sugestões foram aceitas apenas quando o autor entendia o código proposto.
Todas foram revisadas, várias foram modificadas e nenhuma foi incorporada sem
teste automatizado correspondente.

\textbf{O que não foi feito com a ferramenta.} A especificação de requisitos,
o projeto da arquitetura, o modelo de dados, a regra de tratamento de
conflitos de reserva (\cref{cod:conflito}) e o texto deste trabalho foram
elaborados sem o uso de ferramentas de \acs{ia}.

\textbf{Configuração.} A extensão foi usada com as configurações padrão,
entre agosto e outubro de 2026, com a opção de compartilhamento de trechos de
código para treinamento desativada. Nenhum dado pessoal, credencial ou dado
dos participantes da avaliação foi fornecido à ferramenta.

% ==============================================================================
% Apêndice B: Planejamento do projeto
% ==============================================================================
\chapter{Planejamento do projeto}
\label{ap:planejamento}

\begin{instrucao}
\orientacao{Planejamento, viabilidade e riscos}
As competências do TCC incluem desenvolver um cronograma, rastrear o progresso
e gerenciar os riscos do projeto. Este apêndice mostra uma forma de registrar
esses elementos. Mantenha o cronograma realizado (e não apenas o planejado) e
comente no texto os desvios relevantes. Apêndices são material elaborado pelo
próprio autor e devem ser citados no texto pelo menos uma vez. Material de
outros autores vai em anexo.

Pela NBR 14724:2024, se um apêndice citar obras, as referências dessas obras
ficam no próprio apêndice, e não na lista geral de Referências.
\end{instrucao}

% ------------------------------------------------------------------------------
\section{Cronograma}

O \cref{qua:cronograma} apresenta o cronograma realizado, organizado nos seis
ciclos de desenvolvimento descritos na \cref{sec:introducao:metodo}. A
avaliação com usuários, prevista para o quinto ciclo, foi adiada para o sexto
porque a implantação em homologação levou mais tempo que o planejado.

\newcommand{\feito}{\cellcolor{gray!45}}
\begin{quadro}[htbp]
  \centering
  \caption{Cronograma realizado por ciclo de desenvolvimento}
  \label{qua:cronograma}
  \small
  \begin{tabular}{|l|c|c|c|c|c|c|}
    \hline
    \textbf{Atividade} & \textbf{C1} & \textbf{C2} & \textbf{C3} &
      \textbf{C4} & \textbf{C5} & \textbf{C6} \\
    \hline
    Levantamento de requisitos         & \feito & \feito &        &        &        &        \\
    \hline
    Protótipos e validação             &        & \feito &        &        &        &        \\
    \hline
    Projeto da arquitetura e dos dados &        & \feito & \feito &        &        &        \\
    \hline
    Implementação (RF01 a RF05)        &        &        & \feito & \feito &        &        \\
    \hline
    Implementação (RF06 a RF08)        &        &        &        & \feito & \feito &        \\
    \hline
    Testes automatizados               &        &        & \feito & \feito & \feito &        \\
    \hline
    Implantação em homologação         &        &        &        &        & \feito &        \\
    \hline
    Avaliação com usuários             &        &        &        &        &        & \feito \\
    \hline
    Escrita do texto                   & \feito &        &        & \feito & \feito & \feito \\
    \hline
  \end{tabular}
  \fonte{Elaborado pelo próprio autor. Cada ciclo teve duas semanas.}
\end{quadro}

% ------------------------------------------------------------------------------
\section{Análise de riscos}

O \cref{qua:riscos} apresenta os principais riscos identificados no início do
projeto, com a probabilidade, o impacto e as ações previstas. O risco R3 se
concretizou no quinto ciclo e a ação de contingência foi aplicada.

\begin{quadro}[htbp]
  \centering
  \caption{Riscos do projeto e ações de prevenção e contingência}
  \label{qua:riscos}
  \small
  \begin{tabularx}{\textwidth}{|c|>{\raggedright\arraybackslash}X|c|c|>{\raggedright\arraybackslash}X|}
    \hline
    \textbf{Código} & \textbf{Risco} & \textbf{Prob.} & \textbf{Impacto} &
      \textbf{Prevenção e contingência} \\
    \hline
    R1 & Pouca disponibilidade dos usuários para entrevistas & Média & Alto &
      Agendar as entrevistas no início do semestre. Usar questionário
      \emph{on-line} se não houver agenda. \\
    \hline
    R2 & Estimativas otimistas de esforço & Alta & Médio &
      Priorizar as histórias com os usuários. Adiar as de baixa prioridade. \\
    \hline
    R3 & Atraso na liberação do servidor de homologação & Média & Alto &
      Solicitar o servidor no segundo ciclo. Implantar em serviço gratuito de
      nuvem como alternativa temporária. \\
    \hline
    R4 & Falta de familiaridade com o arcabouço & Baixa & Médio &
      Fazer um protótipo técnico no primeiro ciclo. \\
    \hline
  \end{tabularx}
  \fonteautor
\end{quadro}

\end{apendicesenv}

% Anexos: material de outros autores (por exemplo, um documento de terceiros
% que o trabalho analisa). Descomente e crie o arquivo se precisar.
% \begin{anexosenv}
% \input{tecnologico/anexos/a-nome-do-anexo}
% \end{anexosenv}

\end{document}
